\documentclass[11pt]{amsart}
\IfFileExists{lmodern.sty}%
  {\usepackage[T1]{fontenc}\usepackage{lmodern}\usepackage{microtype}}%
  {\usepackage[expansion=false]{microtype}}
\usepackage{amssymb,mathtools}
\usepackage{booktabs,array}
\usepackage{enumitem}
\emergencystretch=2.5em
\usepackage[colorlinks=true,linkcolor=blue,citecolor=blue,urlcolor=blue]{hyperref}

% ---------------------------------------------------------------------------
% Location of the Lean formalization.
\newcommand{\FormalizationURL}{https://github.com/slitsey/catalan-irrationality-measure}
\newcommand{\FormalizationLink}{\url{\FormalizationURL}}
% ---------------------------------------------------------------------------

\newtheorem{theorem}{Theorem}[section]
\newtheorem{lemma}[theorem]{Lemma}
\newtheorem{proposition}[theorem]{Proposition}
\newtheorem{corollary}[theorem]{Corollary}
\theoremstyle{definition}
\newtheorem{definition}[theorem]{Definition}
\theoremstyle{remark}
\newtheorem{remark}[theorem]{Remark}
\numberwithin{equation}{section}
\setcounter{tocdepth}{1}

\newcommand{\Z}{\mathbb Z}
\newcommand{\Q}{\mathbb Q}
\newcommand{\R}{\mathbb R}
\newcommand{\C}{\mathbb C}
\newcommand{\F}{\mathbb F}
\newcommand{\N}{\mathbb N}
\newcommand{\den}{\operatorname{den}}
\newcommand{\rank}{\operatorname{rank}}
\newcommand{\corank}{\operatorname{corank}}
\newcommand{\lean}[1]{\nolinkurl{#1}}
\renewcommand{\Re}{\operatorname{Re}}
\renewcommand{\Im}{\operatorname{Im}}
\newcommand{\Lw}{\mathsf L}          % Loewner matrix
\newcommand{\cO}{\mathcal O}

\begin{document}

\title{Irrationality measures for $G\pm5\pi^2/48$ and $L(2,\chi_{-3})$}

\author{Sean Litsey}
\email{sean.litsey.math@gmail.com}

\date{October 19, 2026}

\subjclass[2020]{Primary 11J82; Secondary 11J72, 15A15, 31A15, 68V20}

\keywords{Catalan's constant, $L(2,\chi_{-3})$, irrationality exponent, L\"owner matrices,
Cauchy determinants, Andr\'eief identity, $p$-adic valuations, Jacobi matrices, Darboux transformation,
formal verification, Lean}

\begin{abstract}
Building on OpenAI's determinant method for the irrationality of Catalan's constant
$G=\sum_{j\ge0}(-1)^j/(2j+1)^2$, we prove
\[
\begin{gathered}
 \mu\Bigl(G-\frac{5\pi^2}{48}\Bigr)\le 34.74,\qquad
 \mu\Bigl(G+\frac{5\pi^2}{48}\Bigr)\le 34.74,\\
 \mu\bigl(L(2,\chi_{-3})\bigr)\le 18.50 ,
\end{gathered}
\]
where $\mu$ is the irrationality exponent. We know of no earlier irrationality measure for $G\pm5\pi^2/48$.
For $L(2,\chi_{-3})$, which Calegari, Dimitrov and Tang proved irrational, the previous best bound we know of is their
effective bound $\mu\le24781$.

The approximations come from determinants of \emph{rectangular} L\"owner matrices of digamma differences at integer nodes.
We bound their denominators with $p$-adic assignment problems and rank deficits modulo $p$. We prove nonvanishing
modulo a prime by reducing it to bordered Krylov determinants, which we evaluate through Jacobi matrices and
Darboux transformations. For $L(2,\chi_{-3})$, a finite-field interpolation argument bounds the corank by three;
removing at most three common low nodes gives nonvanishing at a fixed scale throughout a prime interval. We bound their size with Andr\'eief's identity and a certified lower bound for a
discrete logarithmic energy.

Every result is formalized in Lean~4, with no unproved steps and only the standard axioms.
\end{abstract}

\maketitle
\tableofcontents

%% =====================================================================
\section{Introduction}\label{sec:intro}

OpenAI \cite{OpenAI2026} recently proved that Catalan's constant
\[
 G=\sum_{j=0}^\infty\frac{(-1)^j}{(2j+1)^2}=0.91596\,55941\ldots
\]
is irrational, and formalized the proof in Lean. (An earlier claimed proof
\cite{Sun2026} was found to contain a gap \cite{Wachs2026}.) Their argument, and the
irrationality-measure results built on it \cite{Jha2026,Carreira2026,Litsey2026}, use a
determinant of large size whose denominators and decay are balanced very closely.

In the present paper we study a different, much smaller, family of determinants: rectangular
L\"owner matrices of digamma differences, attached to the numbers
\[
\begin{gathered}
 \lambda=\frac G2-\frac{5\pi^2}{96},\qquad \lambda'=\frac G2+\frac{5\pi^2}{96},\qquad
 \frac14L(2,\chi_{-3}),\\
 L(2,\chi_{-3})=\sum_{k\ge0}\Bigl(\frac1{(3k+1)^2}-\frac1{(3k+2)^2}\Bigr)=0.78130\,24128\ldots
\end{gathered}
\]
Thus $G-5\pi^2/48=2\lambda$, $G+5\pi^2/48=2\lambda'$, and $\lambda+\lambda'=G$.

Recall that the \emph{irrationality exponent} $\mu(\xi)$ of a real number $\xi$ is the
infimum of the $\mu$ such that $|\xi-a/b|>b^{-\mu}$ for all integers $a$ and all
sufficiently large $b$.

\begin{theorem}\label{thm:main}
For every $\varepsilon>0$ there is $q_0$ such that for all coprime integers $p$, $q$ with $q\ge q_0$,
\[
\begin{gathered}
 \Bigl|G-\frac{5\pi^2}{48}-\frac pq\Bigr|>q^{-34.74-\varepsilon},\qquad
 \Bigl|G+\frac{5\pi^2}{48}-\frac pq\Bigr|>q^{-34.74-\varepsilon},\\
 \Bigl|L(2,\chi_{-3})-\frac pq\Bigr|>q^{-18.50-\varepsilon}.
\end{gathered}
\]
In particular $G\pm5\pi^2/48$ are irrational, $\mu(G\pm5\pi^2/48)\le34.74$ and
$\mu(L(2,\chi_{-3}))\le18.50$.
\end{theorem}

The computed values of the exponents are $34.738$ and $18.498632$. Theorem~\ref{thm:main}, together with every
result in this paper, has been verified in the Lean~4 proof assistant, with no unproved steps and only the
standard axioms; see Section~\ref{sec:lean}. The formalization is available at \FormalizationLink. The use of AI
in this work is disclosed at the end of the paper.

\subsection{Previous results}
For earlier work on the arithmetic of $G$ and related values see
\cite{RivoalZudilin2003,Zudilin2003,Zudilin02Cheb,KrattRivoal08,KrattZudilin19} and the survey
\cite{Zudilin19Survey}; for recent rational approximations to $G$ see Eskandari \cite{Eskandari2026}. Finite
irrationality exponents for $G$ itself were obtained after OpenAI's proof by Jha \cite{Jha2026}, Carreira
\cite{Carreira2026} and the author \cite{Litsey2026}. Our searches found no previous irrationality result, and no irrationality measure, for $G-5\pi^2/48$ or $G+5\pi^2/48$.
Eskandari, Murty and Nemoto \cite{EMN2025} construct a $2$-dimensional mixed motive over $\Q$ with period $G$ and
use it to obtain linear forms in $1$ and $G$. Eskandari \cite{Eskandari2026} improves the approximation exponent of these
Rivoal--Zudilin type forms from $0.52$ to $0.62$ (a preprint; explicit approximations, not a lower bound for all
rationals). Sun's preprint \cite{Sun2026} claimed a proof of the irrationality of $G$; Wachs \cite{Wachs2026} reports that
its main estimate is not justified.

Calegari, Dimitrov and Tang proved that $1$, $\zeta(2)$ and $L(2,\chi_{-3})$ are linearly independent over
$\Q$ \cite{CDT2024}, by a new arithmetic holonomy bound applied to a construction of Zagier. In
\cite{CDT2025} they made holonomy bounds effective; Theorem~5.1 there gives $\mu(L(2,\chi_{-3}))\le24781$.
Theorem~\ref{thm:main} improves this to $18.50$.

For comparison, among classical constants, the bound $\mu(\log2)<3.5746$ of Marcovecchio \cite{Marcovecchio09}
and the Rhin--Viola bounds for $\zeta(2)$ and $\zeta(3)$ \cite{RhinViola96,RhinViola01} (after Ap\'ery and Beukers
\cite{Beukers79}) have been joined by recent preprints. OpenAI proved $\mu(\pi)=2$, with a Lean formalization \cite{OpenAIPi2026}, improving
the bound $\mu(\pi)<7.1033$ of Zeilberger and Zudilin \cite{ZeilbergerZudilin20}. Liu, Jiang and Zhang proved $\mu(\log a)=2$ for
every rational $a>0$, $a\ne1$, and for nonzero rational combinations of such logarithms \cite{LiuJiangZhang26}. Giraudin
\cite{Giraudin26} lowered the bound for $\mu(\zeta(2))$ to below $5.0193784$. Fauzan \cite{Fauzan2026} announced a proof that
$\zeta(5)$ is irrational, using rationally normalized Hankel determinants, and Lai
\cite{Lai2026} combined his method with simultaneous approximations of Zudilin and of Dauguet and Zudilin to show that
$1$, $\zeta(3)$ and $\zeta(4)$ are linearly independent over $\Q$.

\emph{Related work on determinants.} Irrationality proofs built on determinants of large size have a precedent
in Zudilin's determinantal criterion with Hankel determinants \cite{Zudilin}; the preprint of Liu, Jiang
and Zhang on logarithms \cite{LiuJiangZhang26} uses a separated-weight interpolation determinant, and the arguments of
Fauzan and Lai \cite{Fauzan2026,Lai2026} use Hankel determinants. Brown \cite{Brown26} proposes a multivariate
transfinite-diameter test for Mellin-type integrals, applied to $\zeta(2)$. The framework
we use, determinants whose decay is controlled by Cauchy determinants and Andr\'eief's identity and whose
denominators are controlled prime by prime, is OpenAI's \cite{OpenAI2026}. Our determinants are L\"owner
(divided-difference) matrices of a digamma difference, a matrix-valued analogue of Pad\'e and Hermite--Pad\'e
approximation \cite{NikishinSorokin91}. Formal proofs of irrationality results in Lean also include the
irrationality of $\zeta(3)$ by Liu, Zhang and Zhi \cite{LZZ25}.

\subsection{Outline of the proof}
Section~\ref{sec:overview} explains the construction and the new ideas without technicalities, and states the
general deduction. Sections~\ref{sec:construction}--\ref{sec:decay} give the construction and the three estimates
(denominator, nonvanishing, decay); Section~\ref{sec:proof} deduces Theorem~\ref{thm:main}.
Section~\ref{sec:lean} describes the formalization, and the disclosure of AI use follows it.

\subsection{Notation and conventions}
$\psi=\Gamma'/\Gamma$ is the digamma function; $v_p$ is the $p$-adic valuation; for $P\in\Q[x]$,
$\den P$ is the least positive integer $d$ with $dP\in\Z[x]$; $\vartheta(x;M,a)=\sum_{p\le x,\ p\equiv a\ (M)}\log p$;
$(x)_n=x(x+1)\cdots(x+n-1)$; $[0,r)=\{0,1,\dots,r-1\}$. For finite sets $I,J$ of integers and a matrix $A$
indexed by integers, $A[I,J]$ is the submatrix with rows $I$ and columns $J$.

%% =====================================================================
\section{Overview of the method and of the new ideas}\label{sec:overview}

\subsection{L\"owner matrices of digamma differences}
For a function $F$ analytic near distinct points $x_0,\dots,x_{h-1}$ the \emph{L\"owner matrix}
\cite{Loewner34} is the matrix of divided differences
\[
 \Lw_F[a,b]=\frac{F(x_b)-F(x_a)}{x_b-x_a}\ (a\ne b),\qquad \Lw_F[a,a]=F'(x_a).
\]
It is a Cauchy-like structured matrix, and Hermite's formula gives
$\Lw_F=\frac1{2\pi i}\oint F(z)\,v(z)v(z)^{\top}dz$ with $v_a(z)=1/(z-x_a)$, so that a representation of
$F$ as a Stieltjes or contour integral exhibits $\Lw_F$ as a Gram matrix of Cauchy kernels.

We take three digamma differences, with integer nodes:
\begin{equation}\label{eq:frames}
\begin{gathered}
 F_\lambda(z)=\psi(z+1)-\psi(z+\tfrac34),\quad \omega=\tfrac1{16},\quad \xi=\lambda:=\tfrac G2-\tfrac{5\pi^2}{96},\\
 F_{\lambda'}(z)=\psi(z+1)-\psi(z+\tfrac14),\quad \omega=-\tfrac1{16},\quad \xi=\lambda':=\tfrac G2+\tfrac{5\pi^2}{96},\\
 F_{\chi}(z)=-\tfrac13\psi(z+\tfrac13)+\tfrac12\psi(z+\tfrac12)-\tfrac16\psi(z+1),\\
 \omega=-\tfrac16,\quad \xi=\tfrac14L(2,\chi_{-3}).
\end{gathered}
\end{equation}
In each case $\xi=\omega\,F'(0)$ (this is an evaluation of $\psi'$ at $\frac13,\frac12,\frac14,\frac34,1$), while
$F(b)-F(a)$ and $F'(a)-F'(0)$ are rational for integers $a,b\ge0$. Hence, for any finite sets of
nonnegative integer nodes $R$ (rows) and $Q$ (columns),
\begin{equation}\label{eq:shift}
 \omega\,\Lw_F[R,Q]=M[R,Q]+\xi\,E_{R,Q},\qquad E_{R,Q}[a,b]=\mathbf 1[a=b],
\end{equation}
with an explicit rational symmetric kernel $M$ on $\Z_{\ge0}^2$ (\S\ref{sec:construction}). Replacing $\xi$ by
a variable $x$ gives a matrix polynomial $\Phi(x)=M[R,Q]+xE_{R,Q}$, whose determinant $P(x)$ is a polynomial
with rational coefficients of degree at most $|R\cap Q|$, and $P(\xi)=\det(\omega\,\Lw_F[R,Q])$. If
$|\xi-p/q|$ is very small, $P(p/q)$ is very close to $P(\xi)$, which is tiny; on the other hand
$q^{|R\cap Q|}\den(P)\,P(p/q)$ is an integer. If it is not zero we get a contradiction as soon as
\[
 \underbrace{\log\den P}_{\text{content}}+|R\cap Q|\log q\ <\ \underbrace{-\log|P(p/q)|}_{\text{decay}} .
\]
Three estimates are therefore needed: an upper bound for the \emph{content} $\log\den P$, an upper bound
for $|P(x)|$ near $\xi$ (the \emph{decay}), and \emph{nonvanishing} of $P(p/q)$ for a size that can be
chosen in terms of $q$. We use rectangular geometries, with $R\neq Q$, and all three estimates are new for these constants.

\subsection{New idea 1: rectangular (hybrid) frames}
We take a long common block $C$ of $r$ nodes, which carries the variable $x$, and two short interleaved
blocks of $m\approx r/8$ nodes,
\[
\begin{gathered}
 C=[0,r),\qquad X'=\{r,r+2,\dots,r+2m-2\}\ \text{(rows only)},\\ Y'=\{r+1,r+3,\dots,r+2m-1\}\ \text{(columns only)},
\end{gathered}
\]
so that $R=C\sqcup X'$ and $Q=C\sqcup Y'$ (a second ``member'' adds one more common node $r+2m$). The
determinant $P(x)=\det\Phi(x)$ still has degree $\le r$ (or $r+1$), so the integrality loss
$|R\cap Q|\log q$ is that of $r$ nodes, but the matrix has size $r+m$ and its decay is that of a
configuration of $(1+\alpha)r$ nodes, $\alpha=m/r$. The content grows too, but less: the extra rows and
columns never meet the diagonal. With $\alpha=\frac18$ the content and decay constants become (for $\lambda$)
\[
\begin{gathered}
 \Lambda=3.68274\quad(\text{square frame: }313/105=2.98095),\\ D=3.9028\quad(\text{square frame: }3.1114),
\end{gathered}
\]
and the margin $\eta=D-\Lambda$ rises from $0.1304$ to $0.2201$. The interleaving of $X'$ and $Y'$ is
essential: it keeps the two Cauchy determinants in Andr\'eief's formula comparable (\S\ref{sec:decay}) and
it gives the reduction modulo $\ell$ the bordered structure of \S\ref{sec:nonvanishing}.

\subsection{New idea 2: content by $p$-adic assignment problems and rank deficits}
For a prime $p$ let $v$ be the Gauss valuation on $\Q[x]$ (the minimum of the $p$-adic valuations of the
coefficients); it is multiplicative. Expanding $\det\Phi(x)$ by Leibniz, $-v(\det\Phi)$ is at most the
value $\tau_p$ of the \emph{assignment problem} with weights $W_{ab}=\max(0,-v_p(M[a,b]))$, i.e. the
maximum of $\sum_aW_{a\sigma(a)}$ over bijections $\sigma:R\to Q$. By LP duality, any \emph{potentials}
$u_a,v_b$ with $W_{ab}\le u_a+v_b$ certify $\tau_p\le\sum u_a+\sum v_b$. We give explicit potentials for
every prime power $p^e$ (they depend only on $p^e/r$ and on $p\bmod 4$ or $p\bmod 6$), and sum over primes
with the prime number theorem in arithmetic progressions. This gives a ``tropical'' content
$\Lambda_{\rm trop}=3.74993$ for $\lambda$.

The tropical bound is not sharp. If $(u,v)$ is an optimal dual solution, the scaled matrix
$B=p^{u}\Phi p^{v}$ has entries in the valuation ring $\cO$ of $v$, and
$v(\det\Phi)=-\tau_p+v(\det B)\ge-\tau_p+\corank(B\bmod p)$. So every unit of corank of the reduction of
$B$ modulo $p$ over $\F_p(x)$ saves one factor of $p$. We compute this corank in several ranges of primes,
using the block structure of the reduction (residue classes modulo $p$), Schur complements, the identity
$H_{p-1-n}\equiv H_n\pmod p$ for harmonic numbers, an antisymmetry of Cauchy-type sums, and the
Frobenius--K\"onig theorem. The savings lower the content from $3.74993$ to $3.70343$; with the decay constant below this
lowers the exponent from about $50.0$ to $38.35$.
The uniform rescaled-core argument of Section~\ref{sec:uniform-deficit} adds the missing
interval savings without changing the determinant family, giving content $3.68274$ and exponent $34.74$.

\subsection{New idea 3: nonvanishing through a commuting Jacobi matrix}
Fix a prime $\ell$ in a suitable residue class with $\ell\nmid q$. After scaling by $\ell$ the rows and columns
of nodes beyond a threshold $K\approx\ell/4$, the matrix $\Phi(x)$ reduces modulo $\ell$ to a
\emph{bordered Cauchy--Gram} matrix:
\[
\begin{gathered}
 \ell^{2h}\det\Phi(x)\equiv\text{unit}\times\det\begin{pmatrix}yI_K+N_S&W_{Y}\\ W_X^{\top}&0\end{pmatrix}\pmod\ell,
 \\ y=\text{affine function of }x,
\end{gathered}
\]
where $N_S=\sum_{j=1}^{S}c_jc_j^{\top}$ is a sum of rank-one Gram terms of Cauchy vectors
$c_j=(1/(a+j))_{0\le a<K}$ and $W_X,W_Y$ consist of further Cauchy vectors. The two members of the pair differ by
one Gram term. Their reductions cannot vanish simultaneously unless a \emph{bordered Krylov determinant}
\[
 k=\det\bigl[W_Y\,\big|\,v,\ Av,\ \dots,\ A^{K-m-1}v\bigr],\qquad A=\Pi N_S,\quad v=\Pi c_S,
\]
vanishes modulo $\ell$, where $\Pi$ is the projection onto $\ker W_X^{\top}$ along the span of $W_Y$. We prove
that, with the nodes shifted by a parameter $\theta$, $k(\theta)$ is a constant prime to $\ell$ times a product
of powers of $\theta+c$ with $1\le c<\ell$. The proof is algebraic and uses a structure familiar from
orthogonal polynomials: there is an explicit \emph{Jacobi (tridiagonal) matrix} $\tilde T$ that commutes with
$\Pi$ and with $A$ and acts on the Cauchy vectors by a three-term recurrence (a bispectral situation in
the sense of Gr\"unbaum and Haine \cite{GruenbaumHaine97}). Its LU factorisation, reversed to UL, is
(up to a shift and a diagonal conjugation) the Jacobi matrix of the next smaller instance: a discrete
Darboux, or Christoffel, transformation \cite{SpiridonovZhedanov97,GruenbaumYakimov02}. Krylov determinants of
commuting pairs satisfy a cross-ratio identity, chains of Cauchy vectors give products of Cauchy
determinants, and the remaining Gram determinants are evaluated with Cauchy--Binet, Hankel determinants of
moments, Newton's divided differences and the Chu--Vandermonde identity (which also proves the
Pfaff--Saalsch\"utz-type identity needed). One auxiliary determinant has no product formula and is shown to be nonzero
by positivity. An induction on the size closes. We never need the exact exponents: it suffices to know that the
determinant is an $\ell$-adic unit times powers of $\theta+c$ with $1\le c<\ell$.

Because the pair of members is available for every $S$ in a range, the prime $\ell$ can be chosen
\emph{after} the size $r$, anywhere in an interval $[\frac32r,4r]$. This removes the cost of locating $\ell$
from the exponent (for $\lambda$ and $\lambda'$), and the final exponent is simply $c/\eta$.

\subsection{New idea 4: decay through an energy certificate}
Moving Hermite's contour to two rays in the left half-plane writes $\omega\,\Lw_F[R,Q]$ as a Gram matrix
$\int v_a v_b\,dm$ of a complex measure $m$ (finitely many signed atoms at poles of $F$ plus an exponentially
small density on the rays). Andr\'eief's identity with two families of functions (rows and columns) gives
\[
 |\det\Lw_F[R,Q]|\le\frac1{n!}\int\bigl|\mathsf C_R(z)\bigr|\,\bigl|\mathsf C_Q(z)\bigr|\prod_{i=1}^nd|m|(z_i),
\]
where $\mathsf C_R,\mathsf C_Q$ are Cauchy determinants. Taking logarithms turns the integrand into a discrete
logarithmic energy of the free points $z_i$ against the nodes in an external field. We bound this energy
from below by comparing with a \emph{certificate measure} $\sigma$ (a few hundred point masses with dyadic
coordinates) through a regularised kernel $-\frac12\log(|z-w|^2+\rho^2)$, which is conditionally positive
definite on finite sums by an elementary argument (Frullani's integral and Gaussian sums). Everything left is
a finite list of inequalities between explicit functions on boxes, which we verify in interval arithmetic
inside Lean. The same argument bounds every principal minor along the common block, so the determinant stays
small at every point within $e^{-cr}$ of $\xi$.

\subsection{The deduction}
Write $\eta=D-\Lambda$. The three estimates are combined in \S\ref{sec:proof}; the result is
\[
 \mu(\xi)\le c\Bigl(\frac1\eta+g\,s\Bigr),
\]
where $c$ is the window constant of the decay estimate, $g=\lim r/\ell$ for the sizes at which nonvanishing
is proved, and $s$ is the relative length of the interval in which $\ell$ must be found. For $\lambda$ and
$\lambda'$ the size can be chosen independently of $\ell$, so the term $gs$ disappears and
$\mu\le7.6444/0.22006<34.74$. For $L(2,\chi_{-3})$, the original pair family gives
$6.6510\,(1/0.35954+0.8)<23.82$. Section~\ref{sec:minor-nv} replaces it by a family with at most three
common low-node deletions. This permits a fixed scale across a prime interval and needs only $\ell^2\nmid q$;
the term $gs$ disappears, giving $\mu\le6.6510/(3.7251-10339/3072)<18.50$.

\begin{center}
{\setlength{\tabcolsep}{3pt}
\begin{tabular}{lcccccc}
\toprule
constant & content $\Lambda$ & decay $D$ & window $c$ & $\eta=D-\Lambda$ & $gs$ & exponent\\
\midrule
$G-5\pi^2/48$ & $3.68274$ & $3.9028$ & $7.6444$ & $0.22006$ & $0$ & $34.738$\\
$G+5\pi^2/48$ & $3.68274$ & $3.9028$ & $7.6444$ & $0.22006$ & $0$ & $34.738$\\
$L(2,\chi_{-3})$ & $10339/3072$ & $3.7251$ & $6.6510$ & $0.35954$ & $0$ & $18.498632$\\
\bottomrule
\end{tabular}}
\end{center}

\subsection{Formal verification}
Every step, including the arithmetic of the dual certificates, the closed-form evaluations behind the
nonvanishing, the interval-arithmetic decay certificates and the prime number theorem in progressions
modulo $4$ and $6$, is formalised in Lean~4 with Mathlib \cite{Lean4,Mathlib}. The formal proofs are built
on a shared core (rectangular determinants, Gauss valuations and assignment bounds, bordered Krylov
determinants, energy certificates, the measure deduction) with thin instances for the three functions in
\eqref{eq:frames}. The written proofs below are at the level of a research article; for long
computations we point to the formal proof, which is the basis for the correctness of
Theorem~\ref{thm:main}.

%% =====================================================================
\section{The construction}\label{sec:construction}

\subsection{The kernels}
For the three functions of \eqref{eq:frames} the kernel $M$ of \eqref{eq:shift} is as follows; here
$0\le a<b$ are integers and $M[b,a]=M[a,b]$.
\begin{align*}
 \lambda:&\quad M[a,b]=-\frac1{b-a}\sum_{k=a+1}^{b}\frac1{16k(4k-1)},\\
 &\quad M[a,a]=-\sum_{k=1}^{a}\Bigl(\frac1{16k^2}-\frac1{(4k-1)^2}\Bigr);\\
 \lambda':&\quad M[a,b]=\frac1{b-a}\sum_{k=a+1}^{b}\frac3{16k(4k-3)},\\
 &\quad M[a,a]=-\sum_{k=1}^{a}\Bigl(\frac1{(4k-3)^2}-\frac1{16k^2}\Bigr);\\
 \chi:&\quad M[a,b]=\frac1{6(b-a)}\sum_{k=a+1}^{b}d(k),\\
 &\quad M[a,a]=-\sum_{k=1}^{a}\delta(k),
\end{align*}
with $d(k)=\frac1{3k-2}-\frac1{2k-1}+\frac1{6k}$ and
$\delta(k)=\frac1{2(3k-2)^2}-\frac1{3(2k-1)^2}+\frac1{36k^2}=-\frac16d'(k)$.

\begin{lemma}\label{lem:loewner}
For each of the three functions, $\omega\,\Lw_F[R,Q]=M[R,Q]+\xi E_{R,Q}$ for all finite sets $R,Q\subset\Z_{\ge0}$, with
$\omega$ and $\xi$ as in \eqref{eq:frames}.
\end{lemma}
\begin{proof}
From $\psi(w+1)-\psi(w)=1/w$ one gets $F(k)-F(k-1)$ as an explicit rational function of $k$; for instance
$F_\lambda(k)-F_\lambda(k-1)=\frac1k-\frac1{k-1/4}=-\frac1{k(4k-1)}$, and telescoping gives the off-diagonal
entries. Similarly $\psi'(w+1)-\psi'(w)=-1/w^2$ gives $F'(a)-F'(0)$. Finally $\omega\,F'(0)=\xi$: for $\lambda$,
$\frac1{16}(\psi'(1)-\psi'(\frac34))=\frac1{16}(\frac{\pi^2}6-\pi^2+8G)=\lambda$; for $\lambda'$ use
$\psi'(\frac14)=\pi^2+8G$; for $\chi$ use $\psi'(\frac12)=\frac{\pi^2}2$ and
$\psi'(\frac13)=\frac{2\pi^2}3+\frac92L(2,\chi_{-3})$, which give
$-\frac16F_\chi'(0)=\sum_{k\ge1}\delta(k)=\frac14L(2,\chi_{-3})$.
\end{proof}

Note that $2\lambda=G-5\pi^2/48$, $2\lambda'=G+5\pi^2/48$ and $4\xi_\chi=L(2,\chi_{-3})$, and
$\mu(a\xi+b)=\mu(\xi)$ for rational $a\ne0$, $b$. So Theorem~\ref{thm:main} is a statement about the three
numbers $\xi$.

\subsection{Geometries}
\begin{definition}\label{def:geometry}
Let $r,m\ge0$ and $e\in\{0,1\}$. The \emph{geometry} $\mathcal G(r,m,e)$ consists of the common nodes
$C=[0,r)\cup\{r+2m:\ e=1\}$, the row nodes $X'=\{r+2i:0\le i<m\}$ and the column nodes
$Y'=\{r+2i+1:0\le i<m\}$. Rows are indexed by $R=C\sqcup X'$, columns by $Q=C\sqcup Y'$. The
\emph{member} $e=0$ is $\mathcal G_1(r,m)$, the member $e=1$ is $\mathcal G_2(r,m)$. The matrix of the
geometry is
\[
 \Phi(x)=\Phi_{\mathcal G}(x)=\bigl(M[a,b]+x\,\mathbf 1[a=b\in C]\bigr)_{a\in R,\ b\in Q},
\]
a square matrix of size $|C|+m$, and $P_{\mathcal G}(x)=\det\Phi_{\mathcal G}(x)\in\Q[x]$.
\end{definition}

By Lemma~\ref{lem:loewner}, $\Phi_{\mathcal G}(\xi+t)=\omega\,\Lw_F[R,Q]+t\,E_{R,Q}$ for real $t$. The geometry
with $m=0$, $e=0$ is the square L\"owner matrix of size $r$. We use \emph{admissible} geometries, with
$m\approx r/8$: for $\lambda,\lambda'$ the condition is $r\le8m\le r+8$, and for $L(2,\chi_{-3})$ it is
$r\le8m\le r+9$ (the extra unit is needed to find a pair of members for every prime $\ell\equiv1\pmod6$; see
\S\ref{sec:proof}). All $o(1)$ terms below are uniform over admissible geometries as $r\to\infty$.

\begin{lemma}[integrality]\label{lem:integrality}
$P_{\mathcal G}$ has degree at most $|C|$. If $d_{\mathcal G}=\den P_{\mathcal G}$ and $p/q$ is a fraction in
lowest terms with $P_{\mathcal G}(p/q)\ne0$, then $|d_{\mathcal G}P_{\mathcal G}(p/q)|\ge q^{-|C|}$.
\end{lemma}
\begin{proof}
The variable occurs only in the $|C|$ diagonal entries indexed by $C$, so each Leibniz term has degree at most
$|C|$. Then $q^{|C|}d_{\mathcal G}P_{\mathcal G}(p/q)$ is a nonzero integer.
\end{proof}

\begin{remark}
When the $m\times m$ block $M[X',Y']$ is invertible, the Schur complement gives
$P_{\mathcal G}(x)=\det M[X',Y']\cdot\det(xI_C+M_{\rm S})$ with
$M_{\rm S}=M[C,C]-M[C,Y']M[X',Y']^{-1}M[X',C]$. Thus the hybrid determinant is, up to a constant, the
characteristic polynomial of a rational matrix of size $|C|$, as in the square case. We do not need this
factorisation: Lemma~\ref{lem:integrality} only uses the degree bound.
\end{remark}

%% =====================================================================
\section{Arithmetic: the content}\label{sec:content}

In this section we prove the following.

\begin{theorem}[content]\label{thm:content}
For admissible geometries of scale $r$,
\[
 \log\den P_{\mathcal G}\le(\Lambda+o(1))\,r^2,
\]
where $\Lambda=3.68274$ for $\lambda$ and $\lambda'$, and $\Lambda=10339/3072=3.36556\ldots$ for $L(2,\chi_{-3})$.
\end{theorem}

We give the details for $\lambda$ and indicate the changes for $\lambda'$ and for $L(2,\chi_{-3})$ at the end.

\subsection{Gauss valuations and assignment problems}
Fix a prime $p$. Let $\cO\subset\Q(x)$ be the localisation of $\Z_{(p)}[x]$ at the prime ideal $(p)$; it is a
discrete valuation ring with uniformiser $p$ and residue field $\F_p(x)$, and its valuation $v$ restricted to
$\Q[x]$ is the \emph{Gauss valuation} $v(f)=\min_k v_p(f_k)$, $f=\sum f_kx^k$. By Gauss's lemma $v(fg)=v(f)+v(g)$,
and $v_p(\den f)=\max(0,-v(f))$.

\begin{lemma}[assignment bound and LP duality]\label{lem:trop}
Let $B$ be a square matrix over $\Q[x]$ with rows $I$ and columns $J$, $W_{ab}=\max(0,-v(B_{ab}))$, and let
$\tau(W)=\max_\sigma\sum_{a\in I}W_{a\sigma(a)}$ over bijections $\sigma:I\to J$. Then $v(\det B)\ge-\tau(W)$.
If real numbers $u_a$, $v_b$ satisfy $W_{ab}\le u_a+v_b$ for all $a,b$, then $\tau(W)\le\sum_au_a+\sum_bv_b$.
\end{lemma}
\begin{proof}
Each Leibniz term has valuation $\sum_av(B_{a\sigma(a)})\ge-\sum_aW_{a\sigma(a)}$, and $v$ of a sum is at least
the minimum. The second statement is weak LP duality for the assignment problem:
$\sum_aW_{a\sigma(a)}\le\sum_a(u_a+v_{\sigma(a)})$.
\end{proof}
For $\Phi(x)$ the weight of every entry is $W_{ab}(p)=\max(0,-v_p(M[a,b]))$, since $v(x)=0$. We call $\tau_p$ the
value of this assignment problem; thus $v_p(\den P_{\mathcal G})\le\tau_p$.

\subsection{The level majorant}
Let $p\ge5$ and, for a prime power $q=p^e$, let $S_q=\{k\ge1: q\mid k(4k-1)\}$; it consists of the classes
$k\equiv0$ and $k\equiv k_q\pmod q$, where $k_q=(j_qq+1)/4$ with $j_q\in\{1,3\}$, $j_qq\equiv3\pmod4$. Let
$\beta_q(a)=\#(S_q\cap[1,a])$. The ultrametric inequality applied to the formulas for $M$ gives
\[
\begin{gathered}
 W_{ab}(p)\le\sum_{e\ge1}w^{(p^e)}(a,b),\\
 w^{(q)}(a,b)=\mathbf 1[\beta_q(a)\ne\beta_q(b)]+\mathbf 1[q\mid a-b]\ (a\ne b),\\
 w^{(q)}(c,c)=2\cdot\mathbf 1[c\ge k_q],
\end{gathered}
\]
and $w^{(q)}\in\{0,1,2\}$. Since the maximum of a sum is at most the sum of the maxima,
$\tau_p\le\sum_e\tau(w^{(p^e)})$. The primes $2$ and $3$ and all levels $e\ge2$ contribute $o(r^2)$
in total (each $w^{(q)}$ has assignment value $O(r)$, there are $O(\log r)$ levels for each of $2,3$, and
$\sum_{p^e\le Cr,\,e\ge2}\log p=O(\sqrt r)$).

\subsection{Explicit dual certificates}
Fix an odd prime power $q$, write $k=k_q$, $N=r+2m$, and let
\[
 \delta_A(q)=\min\bigl(2r,\ 2(r-n)+2\min(n,(r-q)_+)\bigr),\qquad n=\min(r,k),
\]
the value of the square block $C\times C$. Let $U=[r-q,k)$ (empty if $r-q\ge k$) be the set of ``low unpaired''
indices; a zone node $z\in X'\cup Y'$ is \emph{low} if $z-q\in U$, and $L$ is the number of low zone nodes.

\begin{proposition}[dual certificates]\label{prop:dual}
For $\mathcal G(r,m,0)$ with $8m\le r$, $\tau(w^{(q)})\le\delta_A(q)+E(q)$ with
\begin{center}
\begin{tabular}{llp{0.33\textwidth}}
\toprule
range & $E(q)$ & potentials\\
\midrule
$q<2m$ & $2m-q\,\|m/q\|$ & by residue class mod $q$\\
$2m\le q\le r$ & $\min(2m+L,\ m+L+|U|)$ & two certificates\\
$r<q<N$ & $2m+(N-q)$ & $\mathbf 1[c\ge k]$ on $C$, $1+\mathbf 1[z\ge q]$ on the zone\\
$q\ge N$, $k\le r$ & $2m$ & $\mathbf 1[c\ge k]$ on $C$, $1$ on the zone\\
$q\ge N$, $k>r$ & $(N-k)_+$ & $0$ on $C$, $\mathbf 1[z\ge k]$ on the zone\\
\bottomrule
\end{tabular}
\end{center}
($\|\cdot\|$ is the distance to the nearest integer.)
\end{proposition}
\begin{proof}[Proof sketch]
In each range one writes down potentials and checks $w\le u+v$ entrywise, using two facts: two distinct indices
not congruent modulo $q$ have $w\le1$, and two zone nodes are never congruent once $q\ge2m$. For example, for
$q\ge N$ and $k\le r$, give the common nodes the potential $\mathbf 1[c\ge k]$ and the zone nodes the potential
$1$: there are no congruences, so $w\le1$ off the diagonal, and the diagonal cells (only at $c\in C$) are covered.
For $q<2m$ the potentials are $\frac12,\frac32,1$ on rows and $\frac32,\frac12,1$ on columns of a residue class
$\varrho$, according as $X'$ has more, fewer or as many nodes in $\varrho$ as $Y'$; since the zone alternates
between $X'$ and $Y'$, the imbalance is at most one per class and the number of unbalanced classes is
$2q\|m/q\|$. The certificates for $2m\le q\le r$ start from the square-frame potential on $C$ and raise the
potential of low zone nodes (first certificate), or move half a unit between rows and columns in the residue
classes of high zone nodes (second certificate). All cases are checked in the formal proof.
\end{proof}

\begin{remark}
The certificates are sharp: the exact assignment value computed by the Hungarian algorithm agrees with
$\delta_A+E$ for all odd $q$ and all tested $(r,m)$. Sharpness is not used.
\end{remark}

\subsection{The tropical content}
Put $\theta=q/r$, $\alpha=m/r$. Proposition~\ref{prop:dual} gives $\tau(w^{(q)})\le r\,f_j(\theta;\alpha)+O(1)$
for an explicit continuous piecewise linear function $f_j$ supported in $[0,4(1+2\alpha)/j]$. Summing over
primes in the two classes modulo $4$ with the prime number theorem in arithmetic progressions,
$\frac1r\sum_{p\equiv a\,(4)}f(p/r)\log p\to\frac12\int_0^\infty f$ for bounded Riemann-integrable $f$ of compact
support (see \cite[Ch.~20]{Davenport} for the prime number theorem in progressions), one finds
\[
\begin{gathered}
 \sum_p\tau_p\log p\le(\Lambda_{\rm trop}(\alpha)+o(1))r^2,\\
 \Lambda_{\rm trop}(\alpha)=\frac{313}{105}+\frac{554}{105}\alpha+\Bigl(\frac{118}{21}+2\log2\Bigr)\alpha^2
\end{gathered}
\]
for $\alpha\le\frac18$ (the term $2\alpha^2\log 2$ comes from the sawtooth $q\|m/q\|$ in the range $q<2m$). At
$\alpha=\frac18$, $\Lambda_{\rm trop}=3.7499347\ldots$; for $\alpha=0$ it is the square-frame value $313/105$.

\subsection{Rank deficits}
The assignment bound is sharp for the matrix of weights, but not for $\det\Phi$: cancellation in the
Leibniz expansion makes $\den P$ smaller. We capture it by the following principle.

\begin{lemma}[rank deficit]\label{lem:rank}
Let $(u,v)$ be integer-valued potentials feasible for the weights of $\Phi$ at $p$, of value $\tau'$. Then
$B=\operatorname{diag}(p^{u})\,\Phi\,\operatorname{diag}(p^{v})$ has entries in $\cO$, and
\[
 v_p(\den P_{\mathcal G})\le\tau'-\corank_{\F_p(x)}(B\bmod p).
\]
\end{lemma}
\begin{proof}
$v(\det\Phi)=-\tau'+v(\det B)$, and over the discrete valuation ring $\cO$ the Smith normal form shows
$v(\det B)\ge\corank(B\bmod p)$. (The formal proof avoids the Smith form: if $B\bmod p$ has rank $\rho$, an
explicit rank witness and a change of basis over $\cO$ exhibit $|I|-\rho$ rows divisible by $p$.)
\end{proof}

Write $\delta_p$ for the corank obtained in this way with an optimal dual (so $v_p(\den P)\le\tau_p-\delta_p$).
We use five families of primes. In all of them $p^2>4N$, so only the level $e=1$ matters; the reduction
$B\bmod p$ has a block structure governed by the singular indices $k\in S_p$ below $N$, which cut $[0,N)$ into
blocks on which the leading coefficients ($\operatorname{lead}_{-1}$ of the partial sums of
$1/(4k(4k-1))$ and $\operatorname{lead}_{-2}$ of the diagonal sums) are constant.

\begin{theorem}[rank deficits for $\lambda$]\label{thm:deficits}
Let $\mathcal G=\mathcal G(r,m,e)$ be admissible and $r$ large. Then $v_p(\den P_{\mathcal G})\le\tau_p-\delta_p+O(1)$,
with $\delta_p$ as follows, for all primes $p$ except $p\in\{139,1009\}$:
\begin{center}
\begin{tabular}{p{0.30\textwidth}ll}
\toprule
primes & $\delta_p$ & mechanism\\
\midrule
$p\equiv3\ (4)$, $\frac58r<p<\frac45r$; $\ p\equiv1\ (4)$, $\frac57r<p<r$ & $m$ & integral Schur complement\\
$r<p<r+m$ & $r+m-p$ & harmonic reflection, rank count\\
$p\equiv3\ (4)$, $\frac45(r+m)\le p<r$ & $\frac{5p+1}4-r-m$ & Frobenius--K\"onig\\
$p\equiv1\ (4)$, $\frac47(r+m)\le p<\frac57r$ & $\frac{7p+1}4-r-m$ & Frobenius--K\"onig\\
$p\equiv3\ (4)$, $\frac5{12}r<p<\frac49r$; $\ p\equiv1\ (4)$, $\frac5{11}r\le p<\frac12r$ & $m$ & block-diagonal leading form\\
\bottomrule
\end{tabular}
\end{center}
\end{theorem}

We sketch the mechanisms; the complete proofs are formalised.

\emph{The integral Schur complement.} Write $\Phi=\begin{psmallmatrix}A&\Phi[C,Y']\\ \Phi[X',C]&\Phi[X',Y']\end{psmallmatrix}$
with $A=\Phi[C,C]$, so $\det\Phi=\det A\cdot\det T$ with the Schur complement
$T=\Phi[X',Y']-\Phi[X',C]A^{-1}\Phi[C,Y']$. At these primes the square block has $-v(\det A)=2r$, and the
claim is that $T$ has entries in $\cO$, whereas the assignment bound allows $v(T_{zy})=-1$. For a zone row $z$
with partner $c_z=z-p\in C$ decompose $\Phi[z,C]=A[c_z,C]+\ell_z$, and similarly for columns; since
$A[c_z,\cdot]A^{-1}=e_{c_z}^{\top}$,
\[
\begin{gathered}
 T_{zy}=\Delta^2M(z,y)-\ell_zA^{-1}\ell_y^{\top},\\
 \Delta^2M(z,y)=M(z,y)-M(z,c_y)-M(c_z,y)+M(c_z,c_y),
\end{gathered}
\]
and $x$ cancels. The leading form of $p^2A$ is block diagonal by residue class (blocks of size $1$ and $2$), with
explicit nonzero determinants (here $p\ne139$ is used: one block has determinant $3^2\cdot5\cdot139/6400$).
Expanding $\ell_zA^{-1}\ell_y^{\top}$ modulo $p^2$, the only surviving terms are
$p\beta\bigl[\frac{\sigma}{4(c_z-c_y)}+\frac{\sigma}{4(c_y-c_z)}\bigr]=0$, an antisymmetry of Cauchy-type
sums. For $p\equiv1\pmod4$ the analogous computation uses a scalar Schur complement: modulo $p$ it equals
$s_0I$ with $s_0=x-\frac1{16}H^{(2)}_{K_p-1}$, $K_p=(3p+1)/4$, which follows from partial fractions and the
reflections $H_{p-1-n}\equiv H_n$, $H^{(2)}_{p-1-n}\equiv-H^{(2)}_n\pmod p$ for harmonic numbers. Here it is
essential that $x$ is an indeterminate: $s_0$ is a unit of $\cO$, although it may vanish at a specialisation
$x\in\Q$.

\emph{Primes just above $r$.} For $r<p<r+m$ the only singular indices below $N$ are $k_p$ and $p$. After
eliminating an invertible block, the Schur complement on $[0,k_p)$ is $s_0I-\frac1{16}EE^{\top}$ with
$E=(1/(u+i))_{u<k_p,\,1\le i\le p-r}$, by the same harmonic reflection. Each zone node $z<p$ contributes
one rank-one term indexed by $i=p-z$, so the correction has rank at most $p-r$ and the corank is at least
$r+m-p$.

\emph{Primes with one singular zone index.} For $p\equiv3\pmod 4$ with $\frac45(r+m)\le p<r$ (and
$p\equiv1\pmod 4$ in the corresponding range) there is one singular index inside the zone. After
elimination the zone$\times$zone block of $B\bmod p$ has a zero sub-block of size $L_X\times L_Y$ with
$L_X+L_Y=L\ge m$ (the low zone nodes), so by the Frobenius--K\"onig theorem its rank is at most $2m-L$, and the
corank is at least $L-m$.

\emph{The low band.} For primes $p$ with $2m<p<r$ that are \emph{clean} (no singular index in the zone) we prove
a general statement: if every residue-class block of the leading form of $p^2A$ is invertible modulo $p$, then
$T$ is integral. The proof splits $p^2A=B_0+pE$ with $B_0$ block diagonal by residue class, inverts by a Neumann
series, and shows that the order-$p$ term vanishes: the translation $c\mapsto c+(y-z)$ identifies the classes of
$z$ and $y$ (cleanness is used here), and what remains is $\sum_{i,j}g_ig_j(L_i-L_j)=0$, an antisymmetric sum.
The original proof for the two listed families used classes of size $2$ or $3$ and explicit rational block determinants,
for example $-350972849/368640000$ (with prime numerator), which are nonzero modulo $p$ except at $p=139$ and
$p=1009$. The uniform argument in Section~\ref{sec:uniform-deficit} proves the needed clean-case bound without these exceptional determinants.

\subsection{Summation}
Each family is a union of intervals $p\in(ar,br)$ in a fixed class, and $\delta_p$ is a linear function of $p$
there. With the prime number theorem in progressions, the total saving is
\[
 \Bigl(\frac{129}{4480}+\frac{59}{4480}+\frac{29}{6336}\Bigr)r^2=(0.0287946+0.0131696+0.0045770)\,r^2
\]
(the three terms come from the first row, the next three rows, and the last row of the table). Hence
\[
 \Lambda_\lambda\le\Lambda_{\rm trop}(\tfrac18)-0.0465412=3.7033935\ldots\le3.70343 .
\]
(The formal proof uses the slightly coarser rational value $3.7034204$, from a truncation of the sawtooth term.)
This is the intermediate content constant before the additional savings of Section~\ref{sec:uniform-deficit},
which complete the improved constant in Theorem~\ref{thm:content}.

\subsection{The second member}
The geometry $\mathcal G_2$ has one extra common node $b=r+2m$. Bordering $\Phi_{\mathcal G_1}$ by one row and one
column changes each per-prime bound by $O(1)$, so the content constant is unchanged; but the rank deficits have
to be re-proved for the bordered matrix, since a deficit is not automatically inherited under bordering. For the
deficits proved by rank witnesses there is a general bordering lemma: a witness for $\Phi_{\mathcal G_1}$
extends, with the border row and column scaled by $p^{w}$, to a witness for $\Phi_{\mathcal G_2}$ at a cost of
$2w=O(1)$. For the low-band deficit, which is proved through integrality of the Schur complement, the border row
and column are rescaled by $p^w$ so that the same core argument applies, with one unit of slack.

\subsection{The mirror constant and $L(2,\chi_{-3})$}
For $\lambda'$ the singular indices are those $k$ with $p\mid k(4k-3)$; the two residue classes modulo $4$ exchange
their roles. The tropical law is the same function $\Lambda_{\rm trop}$, the deficits are the same with the classes
swapped, and $\Lambda_{\lambda'}=\Lambda_\lambda$. The formal proof uses one generic content theorem with the offset
($1$ or $3$) as a parameter.

For $L(2,\chi_{-3})$ the singular indices are those $k$ with $p\mid k(3k-2)(2k-1)$, and the relevant classes of
primes are modulo $6$. Here no rank deficits are used. Two explicit duals suffice: below a threshold $\theta^*$
($\frac{27}{32}$ for $p\equiv1$, $\frac34$ for $p\equiv5\pmod6$) the residue-class potentials of the first row of
Proposition~\ref{prop:dual}, and above $\theta^*$ a certificate in which only row potentials are raised, on nodes
that have a congruent partner of level $0$. The levels (leading coefficients of the partial sums) enter only
through which of them vanish, and this reduces to eight explicit identities. The resulting bound is piecewise
linear in $p/r$, with $3+6$ and $3+5$ pieces in the two classes, and integrating gives
$\Lambda_\chi=\frac{10339}{3072}$. The per-prime bound agrees with the exact assignment value at all primes
tested ($r=256$ and $r=1024$).

%% =====================================================================
\section{A uniform rescaled-core deficit}
\label{sec:uniform-deficit}

We retain the determinant family and the previously established clean-prime
estimates.  A common rescaled-core argument supplies three additional,
disjoint groups of prime windows.  The gain comes from a first-order
cancellation and an elementary assignment dual; no assertion about a Smith
normal form or the optimality of the previous tropical bound is required.
The argument applies to both offsets $o\in\{1,3\}$.

\paragraph{Monic rescaling.}
For the finite certificate take $m=\lfloor r/8\rfloor$; the same leading estimates hold for
$m=r/8+O(1)$, since all endpoint and border losses are bounded per prime. Put $N=r+2m$ and write $C=[0,r)$ for the
common core and $Z=[r,N)$ for the interleaved border zone.  We use primes
$p\ge13$ satisfying $4(N+1)<p^2$.  In the normalized symmetric rational
matrix $R$, whose omitted normalization factors are $p$-adic units, every
entry has valuation at least $-2$.  Replace the polynomial variable $X$
by $x/p^2$ and set
\[
 A=R_{CC}+\frac{x}{p^2}I_r,\qquad Q=p^2A=p^2R_{CC}+xI_r.
\]
The entries of $Q$ are integral for the $p$-adic Gauss valuation.  Its
reduction is block diagonal according to residues of core nodes modulo
$p$.  Every reduced block has the form $M+xI$, so its determinant is
monic and nonzero over $\mathbb F_p(x)$.  Consequently
\[
 v_G(\det Q)=0,\qquad v_G(\det A)=-2r.
\]
This remains valid when the constant matrix $M$ is singular.  The single
Gauss unit $\det Q$ and its adjugate suffice to express all inverse
entries; separate denominators for individual residue blocks are unnecessary.

Let $T$ be the Schur matrix left after eliminating $C$.  If every
assignment in its $m\times m$ zone matrix has total entry price at most
$\tau$, with $v_G(T_{zy})\ge-w_{zy}$, then the determinant expansion gives
\[
 v_G\bigl(F(x/p^2)\bigr)\ge-2r-\tau.
\]
Indeed, for $F(X)=\sum_i c_iX^i$ the left side is
$\min_i(v_p(c_i)-2i)$.  Thus every coefficient satisfies
$v_p(c_i)\ge-2r-\tau+2i\ge-2r-\tau$.
The substitution estimates coefficients of the original polynomial and
introduces no new approximation family.

\paragraph{Translation and cancellation.}
The singular cuts are the integers in the two residue classes
$0$ and $(cp+o)/4$ modulo $p$, where $c=1$ or $3$ is determined by the
prime class.  For a zone node $z$, its core chain consists of
\[
 \{z-ip: i\ge1,\ 0\le z-ip<r\}.
\]
Between consecutive singular cuts and the core-exit cuts $r+ip$, translation
by $y-z$ identifies the chains of $z$ and $y$, their leading node data,
the leading core--zone vectors, and their symmetric leading chain matrices.
For the one-interval construction below, clipping the right endpoint at
$r+p$ controls the core-exit cut.  The formal diagonal $xI$ is preserved
by the same identification.

For completeness, write $Q=B+pE$, where $B$ retains the full
same-residue blocks of $Q$.  Modulo $p^2$,
\[
 Q^{-1}=B^{-1}-pB^{-1}EB^{-1}.
\]
For two distinct residues the order $-2$ Schur term vanishes.  For two
nodes in one translated class, let $g$ be the common leading inverse
vector and $L$ the leading node datum.  Up to one common normalization,
the three order $-1$ correction sums are
\[
 \begin{gathered}\frac\gamma{y-z}\sum_u g_u\bigl(L(z)-L(u)\bigr),\\
 \frac\gamma{y-z}\sum_u g_u\bigl(L(u)-L(z)\bigr),\\
 -\frac\gamma{y-z}\sum_{u,v}g_ug_v\bigl(L(u)-L(v)\bigr).\end{gathered}
\]
The first two cancel termwise, and the last vanishes by interchanging
$u,v$.  The direct order $-1$ entry also vanishes since $L(y)=L(z)$.
This proves integrality within the selected class.  Between distinct
residues the price is at most $1$; equal residues have price at most $2$.
No specialization at $x=0$ or additional spectral hypothesis enters this
calculation.

\paragraph{One interval and a $0/1/2$ dual.}
Choose consecutive singular cut indices $A_4<B_4$, with no permitted
numerator strictly between them, and put
\[
 \begin{gathered} a=\max\left(r,\left\lfloor\frac{A_4p+\epsilon_{A_4}}4\right\rfloor\right),
 \\ b=\min\left(N,r+p,
       \left\lfloor\frac{B_4p+\epsilon_{B_4}}4\right\rfloor\right),
 \\ \epsilon_j=\begin{cases}0&j\equiv0\pmod4,\\o&\text{otherwise}.\end{cases}\end{gathered}
\]
In the high range $p>2m$ the clipping at $r+p$ is redundant.
All nodes of $[a,b)$ lie in one translated class and have distinct
residues.  In the low range $m<p<2m$, use the smaller interval
\[
 J=[\max(a,N-p),b),
\]
interpreting its length as zero when its left endpoint exceeds its right
endpoint.  Since $b\le r+p$, a node in $J$ has no other zone node of the
same residue.  In the high range take $J=[a,b)$.
Assign potential $0$ to every row or column node in $J$, and potential
$1$ to every other zone node.  Two nodes in $J$ have integral Schur
entry, a mixed pair has price at most $1$, and an outside pair has price
at most $2$.  Therefore this is a feasible dual and
\[
 \tau\le2m-|J|.
\]
This elementary certificate replaces class-count and parity calculations.

The second member adds one row and column to the remaining Schur matrix.
Every new incident entry, including its rescaled polynomial corner, has
price at most $2$.  An assignment uses at most two such incidences;
extending the remaining partial assignment gives a total additional
price at most $4$.  Thus both members satisfy the common coefficient
bound
\[
 \operatorname{denexp}_p F_i\le2r+2m-|J|+4\qquad(i=1,2).
\]
The constant $4$ is retained in the integer deficit, rather than silently
discarded before summation.

\paragraph{Baseline and finite profile.}
Let $T_{\rm old}$ be the exponent used for the already banked estimates.
Only on the new profile windows enlarge it to
\[
 T_*=\max(T_{\rm old},B_0),\qquad
 B_0=\begin{cases}2r+3m-p,&m<p<2m,\\2r+m,&p>2m.\end{cases}
\]
The existing affine majorants already cover these values with their
existing constant slack; their leading integral is unchanged.  This
comparison requires no equality between $B_0$ and a tropical optimum.
The resulting nonnegative integer deficits are
\[
 \delta_p=\begin{cases}
 (m+|J|-p-4)_+,&m<p<2m,\\
 (|J|-m-4)_+,&p>2m.
 \end{cases}
\]
Outside the new windows set the new deficit to zero.  Retain all old
E, B, C5, C7, and clean GK estimates: this rescaling is not asserted to
reproduce their sharper savings.

Here is a compact specification of the complete rational certificate.
Write $u=p/r$, $A=A_4/4$, $B=B_4/4$.  In the low range put
\[
 \begin{gathered}a_0=\max(1,Au),\qquad b_0=\min(5/4,Bu,1+u),\\
 f(u)\le\min(1/8-u+b_0-a_0,\ b_0-9/8).\end{gathered}
\]
In the high range the corresponding certificate is
\[
 f(u)\le\min(5/4,Bu)-\max(1,Au)-1/8.
\]
For an affine $f$, each assertion reduces to finitely many affine
endpoint inequalities: nine in the low case and four in the high case.
The certificate contains 37 low rows, 20 nonclean rows below $u=3/5$,
and four missing upper rows.  In addition to the affine inequalities,
it checks consecutive pole numerators, the residue class, nonnegativity,
and pairwise disjointness within each prime class.

Every raw interval $[s,t]$ is replaced once and for all by
\[
 [s+\eta,t-\eta],\qquad \eta=10^{-7}.
\]
This fixed trimming leaves a strict margin at all old-window boundaries.
For either offset the active rows eventually avoid every banked deficit
window; this is a finite rational range check.  In particular there is
no addition of two deficits at the same prime.  The exact integer
implementation also guards the new deficit by the vanishing of the old
one.  The guard holds throughout every trimmed profile window for all
sufficiently large $r$.

Replacing $m$ by $r/8$ and the cuts by $A_4p/4,B_4p/4$ changes interval
endpoints by bounded amounts.  The explicit rounding bounds yield,
uniformly over all rows,
\[
 \delta_p\ge r f(p/r)-8.
\]
Both offsets use the same rational profile.  Its tag $c=1$ selects
$p\equiv4-o\pmod4$, and its tag $c=3$ selects $p\equiv o\pmod4$;
thus the mirror exchanges the two prime classes without changing their
integrated contribution.

\paragraph{Prime summation and the final constants.}
For a row $f(u)=\alpha+\beta u$ on a trimmed interval $[s,t]$, the prime
number theorem in progressions gives the quadratic contribution
\[
 \frac12\int_s^t f(u)\,du
 =\frac12\left(\alpha(t-s)+\frac\beta2(t^2-s^2)\right).
\]
The constant loss $8$ contributes only $O(r)$ to the logarithmically
weighted prime sum.  The finite sum uses the same prime range as the
old denominator bound.  The untrimmed contributions are
{\renewcommand{\arraystretch}{1.8}
\[
\begin{array}{c|c|c}
\text{group}&\text{number of rows}&\text{gain}\\\hline
m<p<2m&37&\frac{14667950664797}{5028641371353600}\simeq0.0029168814\\
\text{nonclean}&20&\frac{160166243}{12415603200}\simeq0.0129003996\\
\text{upper windows}&4&\frac{109}{22400}\simeq0.0048660714
\end{array}
\]
}
The four upper affine profiles, in tag order, are
{\renewcommand{\arraystretch}{1.8}
\[
\begin{array}{c|c|c}
c&u\text{-interval}&f(u)\\\hline
1&(4/5,9/10)&9/8-5u/4\\
1&(3/5,5/8)&2u-9/8\\
3&(5/8,9/14)&9/8-7u/4\\
3&(3/5,5/8)&u/4-1/8.
\end{array}
\]
The 61-row sum after the fixed trimming is exactly
\[
 \Delta_\eta=
 \frac{4420349242147094110541}{213717258282528000000000}
 =0.0206831646525407\ldots.
\]
The trimming loss is less than $1.879\times10^{-7}$.
Starting from the banked content constant
\[
 \begin{aligned}\Lambda_{\rm old}&=\frac{24729491161}{6665218560}-\frac1{448}-\frac1{576}-\frac1{352}
\\ &=\frac{320572813}{86561280}=3.7034204323226274\ldots,\end{aligned}
\]
the new quadratic content constant is at most
\[
 \Lambda_{\rm old}-\Delta_\eta
 =3.6827372676700867\ldots
 \ <\ \frac{184137}{50000}=3.68274.
\]
In the existing measure conversion take $D=3.9028$ and $c=7.6444$.
The rational inequalities
\[
 \frac{c}{D-3.68274}<\frac{1737}{50}=34.74,
 \qquad D-3.68274>0
\]
give $\mu(G\pm5\pi^2/48)\le34.74$ by the existing measure conversion.  This improvement
uses no claimed new gain for $p<m$ and no unproved higher-order Smith
factorization.


\section{Nonvanishing modulo a prime}\label{sec:nonvanishing}

This section follows the formal proof. We treat $\lambda$ in detail; $\lambda'$ and $L(2,\chi_{-3})$ are in
\S\ref{ssec:nvother}. Throughout, $\ell\equiv3\pmod4$ is a prime, $K=(\ell+1)/4$, and $\mathcal G_e=\mathcal G(r,m,e)$,
$e\in\{0,1\}$, is the pair of members of a scale $r$, with
\begin{equation}\label{eq:nvrange}
 1\le m<K\le r,\qquad S:=\ell-r-2m\ \ge\ n:=K-m .
\end{equation}
For fixed $r$ and $r\le8m\le r+8$, \eqref{eq:nvrange} holds for every prime $\ell$ with $\frac43(r+m)+\frac13\le\ell\le4r-1$,
an interval $[\frac32r+O(1),4r]$.

\begin{theorem}[nonvanishing]\label{thm:nv}
Assume \eqref{eq:nvrange} and let $x_0\in\Q$ with $\ell\nmid\den x_0$. Then $P_{\mathcal G_0}(x_0)\ne0$ or
$P_{\mathcal G_1}(x_0)\ne0$.
\end{theorem}

The proof has four parts: a reduction modulo $\ell$ to a pair of bordered Cauchy--Gram determinants
(\S\ref{ssec:red}); a lemma showing that if both vanish, a bordered Krylov determinant vanishes (\S\ref{ssec:primal});
a deformation in a parameter $\theta$ (\S\ref{ssec:deform}); and an algebraic evaluation of the deformed Krylov
determinant, up to factors that are units at $\ell$, by a commuting Jacobi matrix, diagonal steps and Darboux steps
(\S\S\ref{ssec:jacobi}--\ref{ssec:induction}).

\subsection{Reduction to bordered Cauchy--Gram determinants}\label{ssec:red}
Call a node $b$ \emph{low} if $b<K$ and \emph{high} if $K\le b\le\ell-1$; all nodes of $\mathcal G_e$ are at most
$r+2m\le\ell-1$, and all low nodes are common nodes since $K\le r$. For a high node put $j_b=\ell-b\in[1,\ell-K]$,
and for $j\ge1$ let $w_j=(1/(a+j))_{0\le a<K}\in\F_\ell^K$; all denominators are nonzero modulo $\ell$. For a set $J$
let $W_J$ be the $K\times|J|$ matrix with columns $w_j$, $j\in J$. In these coordinates the row nodes $X'$ and the
column nodes $Y'$ have $J_{X'}=S+\{2,4,\dots,2m\}$ and $J_{Y'}=S+\{1,3,\dots,2m-1\}$. The high nodes \emph{absent}
from the geometry are $\{1,\dots,S\}$ for $\mathcal G_0$ and $\{1,\dots,S-1\}$ for $\mathcal G_1$, whose extra node
$r+2m$ has $j=S$. Put
\[
 N=\sum_{j=1}^{S-1}w_jw_j^{\top},\qquad w=w_S .
\]

\begin{lemma}[$\ell$-adic structure]\label{lem:lstruct}
The only $\ell$-singular factor occurring in the entries $M[a,b]$, $a,b\le\ell-1$, is $4K-1=\ell$. Consequently
$M[a,b]$ is $\ell$-integral if $a,b$ are both low, or both high and distinct; for $a$ low and $b$ high,
$\ell M[a,b]\equiv\frac14w_{j_b}(a)$; for $b$ high, $\ell^2M[b,b]\equiv1\pmod\ell$; and on the low block,
$16\,M[\text{low},\text{low}]\equiv W_{[1,\ell-K]}W_{[1,\ell-K]}^{\top}+H^{(2)}_{K-1}I_K\pmod\ell$, where
$H^{(2)}_n=\sum_{i\le n}i^{-2}$.
\end{lemma}
\begin{proof}
$4k-1\in\{\ell,2\ell,3\ell\}$ with $k\le\ell-1$ forces $k=K$. For $a$ low and $b$ high, the sum defining $M[a,b]$
contains the term $k=K$ once, and $16K=4\ell+4$, $b-a\equiv-(a+j_b)$. For the low block, $\frac4{4k-1}\equiv\frac1{k-K}$,
and partial fractions give $\sum_{j=1}^{\ell-K}\frac1{(a+j)(b+j)}$ off the diagonal; on the diagonal one uses
$H^{(2)}_{\ell-K}\equiv-H^{(2)}_{K-1}$.
\end{proof}

\begin{proposition}[reduction]\label{prop:red}
Let $D$ be diagonal with entry $1$ on low and $\ell$ on high rows and columns, and let $h_e$ be the number of high rows of
$\Phi_{\mathcal G_e}$. Then $D\Phi_{\mathcal G_e}(x)D$ is $\ell$-integral, and with $y=16x+H^{(2)}_{K-1}$,
\[
\begin{gathered}
 \ell^{2h_1}P_{\mathcal G_1}(x)\equiv\pm16^{-K}\det\mathcal B(y,N),\\
 \ell^{2h_0}P_{\mathcal G_0}(x)\equiv\pm16^{-K}\det\mathcal B(y,N+ww^{\top})\pmod\ell,
\end{gathered}
\]
where $\mathcal B(y,N')=\begin{pmatrix}yI_K+N'&W_{Y'}\\ W_{X'}^{\top}&0\end{pmatrix}$.
\end{proposition}
\begin{proof}
By Lemma~\ref{lem:lstruct}, modulo $\ell$ the matrix $D\Phi D$ has the block form
$\begin{psmallmatrix}xI+M_{\rm low}&\frac14W_{H}&\frac14W_{Y'}\\\frac14W_{H}^{\top}&I&0\\\frac14W_{X'}^{\top}&0&0\end{psmallmatrix}$,
where rows are low, high common, $X'$, columns are low, high common, $Y'$, and $H$ is the set of $j$-coordinates of the
high common nodes. Take the Schur complement of the identity block. By the low-block formula and the partition of
$[1,\ell-K]$ into $H$, the absent set, $J_{X'}$ and $J_{Y'}$, the upper-left block becomes
$\frac1{16}(yI+N_e+W_{X'}W_{X'}^{\top}+W_{Y'}W_{Y'}^{\top})$, with $N_e$ the Gram matrix of the absent columns. Row and
column operations using the zero lower-right block remove the last two terms, and scaling gives the claim.
\end{proof}

Reduction modulo $\ell$ commutes with evaluation at $x_0\in\Z_{(\ell)}$, and $x\mapsto y$ is the same for both members.
So if $P_{\mathcal G_0}(x_0)=P_{\mathcal G_1}(x_0)=0$, then for $y_0=16x_0+H^{(2)}_{K-1}\in\F_\ell$ both
$\det\mathcal B(y_0,N)$ and $\det\mathcal B(y_0,N+ww^{\top})$ vanish.

\subsection{A pair of singular bordered matrices}\label{ssec:primal}
Let $\Delta=\det(W_{X'}^{\top}W_{Y'})$. If $\Delta\ne0$, let $\Pi=I-W_{Y'}(W_{X'}^{\top}W_{Y'})^{-1}W_{X'}^{\top}$ be the projection
onto $V=\ker W_{X'}^{\top}$ along the column span of $W_{Y'}$, and let $U$ be a $K\times n$ matrix whose columns span $V$.
Define
\[
\begin{gathered}
 k_R=\det\bigl[\,w^{\top}(\Pi N)^kU\,\bigr]_{0\le k<n},\\ k_L=\text{the same with }X'\text{ and }Y'\text{ exchanged}.
\end{gathered}
\]
(Changing $U$ multiplies $k_R$ by a nonzero factor; in $k_L$ the columns of $U$ span $\ker W_{Y'}^{\top}$ instead.)

\begin{lemma}[singular pair]\label{lem:primal}
Let $N'$ be symmetric, $\tau\ne0$, and suppose $\Delta\ne0$. If $\det\mathcal B(y_0,N')=\det\mathcal B(y_0,N'+\tau ww^{\top})=0$
for some $y_0$, then $k_R=0$ for every admissible $U$, or $k_L=0$ for every admissible $U$ (computed with $N'$ in place of $N$).
\end{lemma}
\begin{proof}
Put $\mathcal B_2=\mathcal B(y_0,N')$ and $\hat w=(w,0)$, and consider the bordered matrix
$\mathcal P=\begin{psmallmatrix}\mathcal B_2&\hat w\\-\tau\hat w^{\top}&1\end{psmallmatrix}$. Its Schur complement on the corner $1$
is $\mathcal B_2+\tau\hat w\hat w^{\top}=\mathcal B(y_0,N'+\tau ww^{\top})$, so $\det\mathcal P=0$, and its principal minor $\mathcal B_2$ is
singular. Then $\mathcal B_2$ has a right kernel vector $\hat x$ with $\hat w^{\top}\hat x=0$, or a left kernel vector $\hat x$ with
$\hat x^{\top}\hat w=0$. (If $\rank\mathcal B_2$ is at most its size minus two, its kernel meets $\hat w^{\perp}$. Otherwise
$\operatorname{adj}\mathcal B_2=c\,\hat x\hat y^{\top}$ with $c\ne0$ and $\hat x$, $\hat y$ spanning the right and left kernels, and
$0=\det\mathcal P=\det\mathcal B_2+\tau\,\hat w^{\top}\operatorname{adj}(\mathcal B_2)\hat w=\tau c(\hat w^{\top}\hat x)(\hat y^{\top}\hat w)$.)

Take the right case, $\hat x=(x,c)$. Then $(y_0+N')x+W_{Y'}c=0$, $W_{X'}^{\top}x=0$ and $w^{\top}x=0$; and $x\ne0$, since otherwise
$W_{Y'}c=0$ and $c=0$ because $\Delta\ne0$. Applying $\Pi$ (which kills $W_{Y'}$ and fixes $x\in V$) gives $\Pi N'x=-y_0x$,
hence $w^{\top}(\Pi N')^kx=(-y_0)^kw^{\top}x=0$ for all $k$. Writing $x=Uz$ with $z\neq0$, the matrix
$[w^{\top}(\Pi N')^kU]_k$ annihilates $z$, so $k_R=0$. The left case is the same argument for the transpose,
$\mathcal B(y_0,N')^{\top}=\mathcal B(y_0,N')$ with $X'$ and $Y'$ exchanged, and gives $k_L=0$.
\end{proof}

No condition on $y_0$ is needed; in particular the value $y_0=0$ is not special. By Proposition~\ref{prop:red} and
Lemma~\ref{lem:primal} (with $N'=N$, $\tau=1$), Theorem~\ref{thm:nv} follows from the three statements
\begin{equation}\label{eq:threeunits}
 \Delta\ne0,\qquad k_R\ne0,\qquad k_L\ne0\qquad\text{in }\F_\ell .
\end{equation}

\subsection{Deformation and bordered Krylov determinants}\label{ssec:deform}
We lift to characteristic zero and deform the nodes. Let $\theta$ be a parameter, $c_y=(1/(\theta+a+y))_{0\le a<K}$,
$W_X$, $W_Y$ the matrices with columns $c_y$ for $y\in X=S+\{2,\dots,2m\}$ and $y\in Y=S+\{1,\dots,2m-1\}$,
$N_s=\sum_{y=1}^{s}c_yc_y^{\top}$, $\Pi$ the projection onto $V=\ker W_X^{\top}$ along $\operatorname{span}W_Y$, and
$A_s=\Pi N_s$. For a $K\times K$ matrix $B$ preserving $V$ and $u\in V$, the \emph{bordered Krylov determinant} is
\[
 \kappa(B,u)=\det\bigl[W_Y\,\big|\,u,\ Bu,\ \dots,\ B^{n-1}u\bigr];
\]
it is nonzero exactly when $u$ is a cyclic vector of $B|_V$. Put
\[
 k(\theta)=k_{K,m,S}(\theta)=\kappa\bigl(A_{S-1},\Pi c_S\bigr)=\kappa\bigl(A_S,\Pi c_S\bigr)
\]
(the two agree by (K2) below). At $\theta=0$ all denominators are integers $a+y\le S+K+2m-1=\ell-r+K-1<\ell$.

\begin{lemma}[bridge]\label{lem:bridge}
Let $K=m+n$ with $1\le n\le S$, and let $\theta$ be generic, in the sense that $\theta+z\ne0$ for every integer $z$, with
$\Delta(\theta)\ne0$. Then the row-Krylov determinant $\det[c_S^{\top}(\Pi N_{S-1})^ku_i]_{k,i<n}$, in which $u_1,\dots,u_n$ is the explicit
basis $U$ of $V$ given by Lagrange interpolation at the nodes, equals $\nu(\theta)\,k(\theta)$, where
$\nu(\theta)=C_\nu\prod_{a<K}\prod_{t\in Y}(\theta+a+t)$ and $C_\nu\in\Q^\times$ does not depend on $\theta$. The constant $C_\nu$ is an
$\ell$-adic unit for every prime $\ell>S+K+m$ (all its prime factors are at most $S+K+m$).
\end{lemma}
\begin{proof}[Proof sketch]
Let $U$ and $V'$ be the Lagrange bases of $\ker W_X^{\top}$ and $\ker W_Y^{\top}$, $H=V'^{\top}U$ and $H_c=V'^{\top}N_{S-1}U$; both are
symmetric. Then $\Pi=UH^{-1}V'^{\top}$, hence $\Pi N_{S-1}U=UB$ with $B=H^{-1}H_c$, and $B$ is self-adjoint for $H$, which turns the row-Krylov
matrix into the column-Krylov matrix up to the factors $(c_X/c_Y)^n$ and $\det H$, where $c_X=\Pi_X(S)/\prod_a(\theta+a+S)$ and $c_Y$ is the same with
$\Pi_Y$. Two block-determinant identities,
$\det[W_X\,|\,V']\det[W_Y\,|\,U]=\Delta\det H$ and $\det[W_Y\,|\,C]\det[W_X\,|\,V']=\Delta\det(C^{\top}V')$ with $C=[c_1,\dots,c_n]$,
eliminate $\det H$ and $\Delta$. By Lagrange interpolation $C^{\top}V'$ is a diagonal matrix times a Vandermonde matrix, and $\det[W_Y\,|\,C]$ is a Cauchy
determinant; this gives the stated $\nu$, whose only $\theta$-dependence is $\prod_{a<K}\prod_{t\in Y}(\theta+a+t)$. The operator in the Krylov
determinant of the formalisation is $\Pi N_S\Pi$, which agrees with $A_S$ on $V$ and is related to $A_{S-1}$ by (K2).
\end{proof}

\begin{definition}
Fix $B\ge1$. A function $f$ of $\theta$ is \emph{good} (with bound $B$) if
$f(\theta)=\varkappa\prod_{1\le c\le B}(\theta+c)^{e_c}$ for rational $\theta>0$, with integers $e_c$ and a nonzero rational
$\varkappa$ whose prime factors are at most $B$.
\end{definition}
Good functions with bound $B$ form a group under multiplication. The point of the definition is that a good
function with $B<\ell$ is continuous at $\theta=0$ and its value there is an $\ell$-adic unit.

\begin{theorem}[good form]\label{thm:good}
Assume $1\le m<K$ and $S\ge n=K-m$. Then $k_{K,m,S}$ is good with bound $S+K+2m-1$.
\end{theorem}

Theorem~\ref{thm:good} is all that is needed. The exact exponents $e_c$ are known (they are piecewise quadratic in $c$,
and $\varkappa$ is a ratio of Hankel determinants; this was checked in exact arithmetic in $156$ cases with $K\le9$), but
the proof below never computes them.

\begin{proof}[Proof of Theorem~\ref{thm:nv} from Theorem~\ref{thm:good} and Proposition~\ref{prop:delta}]
By Proposition~\ref{prop:delta} below, $\Delta$ is an $\ell$-unit at $\theta=0$, so $\Delta\ne0$ in $\F_\ell$. The row-Krylov
determinant of Lemma~\ref{lem:bridge} is a continuous function of $\theta$ near $0$ (its entries are rational in
$\theta$ with no pole at $0$, and $\Pi$ is continuous because $\Delta(0)\ne0$). By Lemma~\ref{lem:bridge} and
Theorem~\ref{thm:good} it equals, for rational $\theta\in(0,1)$, a good function with bound $S+K+2m-1<\ell$, which is
continuous at $0$; hence the values at $\theta=0$ agree, and the value is an $\ell$-adic unit. All entries are
$\ell$-integral at $\theta=0$, and reduction modulo $\ell$ of the matrices of \S\ref{ssec:deform} gives the matrices of
\S\ref{ssec:primal} (after reversing the order of the border columns, which does not change $\Pi$). So $k_R\ne0$ in
$\F_\ell$, and the same argument with $X$ and $Y$ exchanged gives $k_L\ne0$. This is \eqref{eq:threeunits}.
\end{proof}

\begin{proposition}[the Cauchy-type Gram determinant]\label{prop:delta}
For $m\le K$ and $\theta\notin\{-1,-2,\dots\}$,
\[
\begin{gathered}
 \det(W_X^{\top}W_Y)=V(X)V(Y)\,\frac{H_{K,m}}{\prod_{i<m}(\theta+S+1+i)_m\,(\theta+S+K+1+i)_m},\\
 H_{K,m}=\det\Bigl[\sum_{a=0}^{K-1}a^{p+q}\Bigr]_{p,q<m},
\end{gathered}
\]
where $V(X)=\prod_{s<s'\in X}(s'-s)$ and $V(Y)$ are Vandermonde products. The Hankel determinant $H_{K,m}$ of the uniform
weight on $\{0,\dots,K-1\}$ is a product of factorials and binomial coefficients with all prime factors less than $2K$.
In particular $\Delta$ is an $\ell$-adic unit at $\theta=0$ when $S+K+2m\le\ell$.
\end{proposition}
\begin{proof}
\emph{Step 1 (Cauchy--Binet).} By Cauchy--Binet and Cauchy's determinant, $\det(W_X^{\top}W_Y)$ is $\pm V(X)V(Y)$ times
$\sum_P V(P)^2\prod_{a\in P}\rho(a)$ over $m$-subsets $P\subseteq\{0,\dots,K-1\}$, where
$\rho(a)=\prod_{s\in X}(\theta+a+s)^{-1}\prod_{t\in Y}(\theta+a+t)^{-1}$. This sum is the Hankel determinant
$\det[\sum_a\rho(a)a^{p+q}]_{p,q<m}$ of the moments of $\rho$.

\emph{Step 2 (basis change and telescoping).} Replacing the monomials $a^p$ and $a^q$ by monic rising factorials
$\varphi_p(a)=(a+\theta+1)_p$ and suitable $\psi_q$ is a unitriangular change of basis, which does not change the
determinant. After it, $\rho(a)\varphi_p(a)\psi_q(a)=1/(a+\theta+p+1)_{2m-p-q}$, and these sums over $a<K$ telescope:
$L\sum_{a<K}1/(c+a)_{L+1}=1/(c)_L-1/(c+K)_L$. Hence each entry is a ratio of Gamma values times a divided difference
$f[u_p,v_q]$ of $f(z)=(z)_K$ at two points $u_p=\theta+p+1$, $v_q=\theta+2m-q$.

\emph{Step 3 (Newton expansion).} We use divided differences only through three facts: the recursion; Newton's
expansion, which peels off the nodes $u_0,u_1,\dots$ one at a time and gives
$f[u_p,v_q]=\sum_{i,j}N_i(u_p)N_j(v_q)\,f[\text{window}]$ with Newton basis polynomials $N_i$; and the values on windows of
consecutive points, $f[s,s+1,\dots,s+n]=\binom Kn(s+n)_{K-n}$ for $f=(z)_K$. Thus $[f[u_p,v_q]]=\mathcal A\,\mathcal D\,\mathcal B^{\top}$ with
$\mathcal A$, $\mathcal B$ triangular, and the determinant reduces to a Pochhammer product times $\det[\binom K{p+q+1}]_{p,q<m}$.

\emph{Step 4 (Chu--Vandermonde).} The identity $\sum_{a=0}^{K-1}\binom ap\binom{K-1-a}q=\binom K{p+q+1}$ (Chu--Vandermonde),
together with another unitriangular basis change, gives $H_{K,m}=\prod_p(-1)^pp!\prod_qq!\,\det[\binom K{p+q+1}]$. No limit
in $\theta$ is used. Finally $H_{K,m}$ is the Hankel determinant of the uniform weight, i.e.\ of the Hahn polynomials with
$\alpha=\beta=0$ \cite[\S9.5]{KLS}; its classical product formula has all factors at most $K+m\le2K-1$.
\end{proof}

\subsection{A commuting Jacobi matrix}\label{ssec:jacobi}
Put $g=S-n+1$ and, for $0\le a<K$,
\begin{align*}
 \sigma(a)&=(a+1)(a+1+\theta)(K-1-a)\frac{\prod_{i=1}^{m}(\theta+S+2i+a)}{\prod_{i=1}^{m-1}(\theta+S+2i+1+a)},\\
 \tau(a)&=(a+1)(a+1+\theta)(K-1-a)\frac{\prod_{i=0}^{m}(\theta+S+2i+1+a)}{\prod_{i=1}^{m}(\theta+S+2i+a)},\\
 \mu(a)&=-(K-1-a)^2(\theta+S+1+a)(\theta+S+2m+1+a).
\end{align*}
Let $\tilde T$ be the $K\times K$ tridiagonal matrix with $\tilde T_{a,a+1}=-\sigma(a)$, $\tilde T_{a+1,a}=-\tau(a)$ and
\[
\begin{gathered}
 \tilde T_{aa}=\mu(a)+\frac{\sigma(a-1)\tau(a-1)}{\mu(a-1)}+g(g+1)a^2+g\,c_0\,a,\\
 c_0=-2(K-1)(\theta+g)-2K^2+2K-m^2 .
\end{gathered}
\]
For $m=0$ this is, up to conjugation, the Jacobi matrix of a family of Hahn polynomials; for $m\ge1$ the products
over $i$ record the border.

\begin{proposition}[bispectral identities]\label{prop:bispectral}
\begin{enumerate}[label=(\roman*),leftmargin=*]
\item \emph{Column chain.} For $1\le y\le S$,
$\tilde Tc_y=l(y)\,c_{y-1}+d(y)\,c_y+r(y)\,c_{y+1}+\sum_{t\in Y}\kappa_t(y)\,c_t$ with $l(y)=-\tau(-\theta-y)$,
$r(y)=-\sigma(-\theta-y-1)$, $d(y)=\tilde T_{aa}|_{a=-\theta-y}$, and explicit corrections $\kappa_t(y)$ given by residues of
$\tau$ and $\sigma$ at the border points. The transpose $\tilde T^{\top}$ satisfies the analogous \emph{row chain} with
corrections in $\operatorname{span}W_X$, with the same coefficients.
\item \emph{Invariance.} $\tilde T$ maps $V=\ker W_X^{\top}$ and $\operatorname{span}W_Y$ into themselves; hence $\tilde T\Pi=\Pi\tilde T$.
\item \emph{Commutation.} $[N_S,\tilde T]$ maps $V$ into $\operatorname{span}W_Y$; hence $\tilde T$ commutes with $A_S$ on $V$.
\end{enumerate}
\end{proposition}
\begin{proof}[Proof sketch]
(i) says that the three-term recurrence of $\tilde T$ in the node variable $a$ is mirrored, under $a\mapsto-\theta-y$, by a
three-term recurrence in the column variable $y$: a bispectral duality in the sense of Gr\"unbaum and Haine
\cite{GruenbaumHaine97}. It is an identity between rational functions of $a$. Conjugating by
$\hat Q(z)=\prod_{t\in Y}(z+t)$ turns $\tilde T$ into a difference operator with polynomial coefficients, and (i) becomes the
statement that a certain polynomial has degree at most $m+2$. Its four top coefficients are computed from the low-order
coefficients of $\prod_i(1+c_iW)$ modulo $W^4$ and vanish identically; the remainder, of degree less than $m$, is
identified by Lagrange interpolation at the border points. (ii) follows from (i) at the border points, where $l$ and $r$
vanish. (iii) follows from (i) and the row chain by telescoping $\sum_{y\le S}\bigl((\tilde Tc_y)c_y^{\top}-c_y(\tilde T^{\top}c_y)^{\top}\bigr)$.
\end{proof}

\subsection{Krylov calculus}
We use five elementary facts about bordered Krylov determinants of matrices preserving $V$.
\begin{enumerate}[label=(K\arabic*),leftmargin=*]
\item \emph{Cross ratio.} If $A$ and $T$ commute on $V$, then $\kappa(A,u)\,\kappa(T,u')=\kappa(A,u')\,\kappa(T,u)$ for all
$u,u'\in V$. (If $T|_V$ is cyclic, $A=p(T)$ on $V$, and $\kappa(A,u)/\kappa(T,u)$ is the determinant of the matrix of the
powers of $p(T)$ in the Krylov basis of $T$, which does not depend on $u$.)
\item \emph{Rank-one update along the start vector.} $\kappa(B+uz^{\top},u)=\kappa(B,u)$.
\item \emph{Chains.} If $Tc_y\equiv l(y)c_{y-1}+d(y)c_y+r(y)c_{y+1}$ modulo $\operatorname{span}W_Y$ for the relevant $y$, then
$\kappa(T,\Pi c_{y_0})=\prod_{k<n}\prod_{j<k}l(y_0-j)\cdot\det[W_Y\,|\,c_{y_0},c_{y_0-1},\dots,c_{y_0-n+1}]$, and similarly upward
with $r$; the Krylov vectors are triangular in the chain.
\item \emph{Cauchy determinants.} $\det[W_Y\,|\,c_{y_1},\dots,c_{y_n}]$ is a Cauchy determinant: a product of differences of
the shifts divided by a product of factors $\theta+a+y$ \cite{Kratt99,Kratt05}.
\item \emph{Transfer.} If $B_2\Phi=\Phi B_1$ for a linear map $\Phi$ carrying $V_1$ into $V_2$, then
$\kappa(B_2,\Phi z)=\det[W_2\,|\,\Phi]\cdot\kappa(B_1,z)$, where $\det[W_2\,|\,\Phi]$ is computed in adapted bases.
\end{enumerate}

\subsection{Proof of Theorem~\ref{thm:good}}\label{ssec:induction}
Write $u(m,n,S;\theta)=k_{m+n,m,S}(\theta)$ and $B=S+n+3m-1=S+K+2m-1$. We prove by induction on $n\ge1$ that
$u(m,n,S)$ is good with bound $B$ for all $S\ge n$.

\emph{Base case $n=1$.} Then $\kappa(A,\Pi c_S)=\det[W_Y\,|\,\Pi c_S]=\det[W_Y\,|\,c_S]$, a Cauchy determinant, which is good.

\emph{Diagonal steps (from $S$ to $S+1$).} Since $c_y(\theta)=c_{y+1}(\theta-1)$, the instances $(\theta,S)$ and
$(\theta-1,S+1)$ have the same border. Apply (K1) to the commuting pairs $(A,\tilde T)$ of both instances, (K2) to pass between
$A_{S-1}$ and $A_S$, and (K3), (K4) to evaluate the four values $\kappa(\tilde T,\Pi c_y)$ that occur as explicit products
(each is nonzero, so $\tilde T$ is cyclic). The result is
\[
 u(m,n,S+1;\theta-1)=u(m,n,S;\theta)\cdot\varkappa_{\rm DS}\prod_{j<n-1}\Bigl(\frac{\theta+K-1-j}{\theta+j}\Bigr)^{2(n-1-j)}
\]
with a constant $\varkappa_{\rm DS}$ whose prime factors are at most $B$. Substituting $\theta+1$ for $\theta$ expresses
$u(m,n,S+1;\theta)$ through $u(m,n,S;\theta+1)$ times factors $\theta+c$ with $1\le c\le K$; so goodness passes from $S$ to $S+1$.
(It is this orientation, downward in $S$ and upward in $\theta$, that keeps all bases $c\ge1$.)

\emph{Darboux steps (from $n-1$ to $n$ at $S=n$).} Consider the instance $(K,m,S)$ with $S=n-1$, i.e.\ $g=0$. After
normalising $\tilde T$ so that it annihilates the kernel vector of $A$ on $V$, it has an LU factorisation
\[
 \tilde T=LU,\qquad U=\operatorname{diag}(\mu)-\operatorname{superdiag}(\sigma),\qquad L=I-\sum_a\frac{\tau(a-1)}{\mu(a-1)}E_{a,a-1},
\]
with explicit pivots $\mu(a)=-(K-1-a)^2(\theta+S+1+a)(\theta+S+2m+1+a)$; note $\mu(K-1)=0$. Reversing the factors gives the
Jacobi matrix of the instance with one node fewer and one more column:
\[
 (UL)_{<K-1}=D^{-1}\bigl(\tilde T'+\beta I\bigr)D,\qquad \tilde T'=\tilde T(K-1,m,S+1),
\]
for an explicit diagonal $D$ and scalar $\beta$. This is the discrete Darboux (Christoffel) transformation of Jacobi
matrices \cite{SpiridonovZhedanov97,Zhedanov97,GruenbaumYakimov02}. The bidiagonal map $X=D\,U$ with its last row removed
is an intertwiner:
\[
 X\tilde T=(\tilde T'+\beta)X,\qquad X\Pi=\Pi'X,\qquad XA=A'X ,
\]
where primes refer to $(K-1,m,S+1)$. The first identity is the reversal of the factors; the other two follow from the
column and row chains of Proposition~\ref{prop:bispectral} at the border points ($Xc_t\in\operatorname{span}W_{Y'}$ for
$t\in Y$, and dually) and, for $A$, from telescoping $XN_S-N'_{S+1}X$ in $y$. By (K5), (K2) and (K1),
\[
 \kappa(A,\Pi c_0)\,\kappa(\tilde T',w')=\det\bigl[W_Y\,\big|\,\Pi c_0\,\big|\,E\bigr]\cdot\kappa(A',w')\cdot
 \prod_{k<n-1}\prod_{j<k}l(-j)\cdot\mathcal G,
\]
where $w'=\Pi'c'_{n}$ is a projected Cauchy vector of the smaller instance, $E$ completes $[W_Y\,|\,\Pi c_0]$ by Cauchy
columns (so the first factor is a Cauchy determinant), and $\mathcal G=\det[R^{\top}\Pi c_{-j}]_{j<n-1}$ with
$R=[c_1,\dots,c_{n-1}]$. Two further inputs are needed.

\begin{enumerate}[label=(\alph*),leftmargin=*]
\item \emph{The Gram determinant $\mathcal G$ (Chu--Vandermonde).} A Schur complement gives
$\mathcal G\cdot\Delta=\det(C_{\mathcal A}^{\top}C_{\mathcal B})$ for two Cauchy matrices whose column sets $\mathcal A$, $\mathcal B$ have $K-1$
elements. By Cauchy--Binet this is a sum over the one omitted node $a$, with
$V(P_a)^2=V_K^2/(a!\,(K-1-a)!)^2$, and the sum is evaluated by
\[
 \sum_{a=0}^{N}\binom Na^2(x+a)_{2N}=\binom{2N}N\bigl((x+N)_N\bigr)^2\qquad(N=K-1).
\]
This is a terminating balanced ${}_3F_2$ identity of Pfaff--Saalsch\"utz type \cite[\S2.2]{AAR99}. We prove it with three
applications of the Chu--Vandermonde identity: the moment identity $\sum_a\binom Na^2a^{\underline j}=N^{\underline j}\binom{2N-j}N$
(falling factorials), the relation $\binom{2N}i\binom{2N-i}N=\binom{2N}N\binom Ni$, and Vandermonde once more for
$(x+2N-1)^{\underline N}$. With Proposition~\ref{prop:delta} this makes $\mathcal G$ explicit.
\item \emph{A nondegeneracy (positivity).} The transfer (K5) needs $\det[W'\,|\,X\Pi R]\neq0$. Its numerator, as a
function of $\theta$, is not a product, so this is proved by positivity rather than by a formula: for rational $\theta>0$
(this is the only range used, and it is proved over $\Q$),
if $W's+X\Pi Rt=0$, then $s=0$ (apply $W_{X'}^{\top}$ and $X\Pi=\Pi'X$), so $y=\Pi Rt$ lies in $V\cap\ker X$; the range of
$X^{\top}$ is contained in the span of $[R\,|\,W_X]$ (the row chain at the border points) and both have dimension $K-1$, so they are
equal and $R^{\top}y=0$, i.e.\ $(R^{\top}\Pi R)t=0$. Finally
$\det(R^{\top}\Pi R)\cdot\Delta=\det([R\,|\,W_X]^{\top}[R\,|\,W_Y])$ by a Schur complement. For two families of $m'$ distinct
shifts $x$, $y$ with $a+x_i>0$ and $a+y_i>0$ for all nodes $a$, Cauchy--Binet gives
\[
 \det(C_x^{\top}C_y)=V(x)V(y)\sum_{P}V(P)^2\prod_{a\in P}\prod_{i}\frac1{(a+x_i)(a+y_i)},
\]
the sum over $m'$-subsets $P$ of the nodes, with $V$ the Vandermonde product. The sum is positive and $V(x)V(y)\neq0$, so the
determinant is nonzero. This applies to $\Delta$ and to $[R\,|\,W_X]^{\top}[R\,|\,W_Y]$, hence $\det(R^{\top}\Pi R)\ne0$ and $t=0$.
\end{enumerate}

Every factor in the displayed identity is now good except the products over $l(-j)=-\tau(-\theta+j)$ and the
Pochhammer denominator of $\mathcal G$, which contain factors $\theta-j-1$ with negative shift; these two sets of factors
coincide and cancel. The values $\kappa(\tilde T',w')$ and $\kappa(\tilde T,\cdot)$ are chain products by (K3) and nonzero.
Combining with the diagonal step at $(K,m,n-1)$, (K1) and (K2), one obtains
\[
 u(m,n,n;\theta-1)=u(m,n-1,n;\theta)\cdot(\text{good factor}).
\]
Since $u(m,n-1,n)$ is good by the induction hypothesis (as $n\ge n-1$), so is $u(m,n,n)$, and the diagonal steps give all
$S\ge n$. All bases stay in $[1,B]$ and all constants have prime factors at most $B$. This proves Theorem~\ref{thm:good}. \qed

\subsection{The mirror constant and $L(2,\chi_{-3})$}\label{ssec:nvother}
\emph{The mirror constant.} For $\lambda'$ one takes primes $\ell\equiv1\pmod4$ and $K=(\ell+3)/4$. The only $\ell$-singular
factor is again $4K-3=\ell$, and the reduction gives the same pair of bordered matrices (with the opposite sign of the
shift). Everything else is unchanged.

\emph{$L(2,\chi_{-3})$: three levels of nodes.} Let $\ell=6t+1$ be prime, $K=(\ell+2)/3=2t+1$, $T=(\ell+1)/2=3t+1$, and
$m_0=3t$. Now there are two $\ell$-singular indices among $1,\dots,\ell-1$, namely $K$ (where $3K-2=\ell$) and $T$
(where $2T-1=\ell$); the factor $k$ in $d(k)$ is a unit. Accordingly the nodes are \emph{low} $[0,K)$, \emph{middle}
$[K,m_0]$ and \emph{top} $[T,\ell-1]$. The geometries are tied to $\ell$: we take $r+2m=m_0$, so all nodes of $\mathcal G_e$
are low or middle, the extra node of the second member is $m_0$, and $K\le r$ is the condition $2m+1\le t$. With the
scaling $1$ on low and $\ell$ on middle and top nodes, a window identity for sliding sums of $1/u$ modulo $\ell$,
\[
 \sum_{a<k\le b}\frac1{k-t_1-1}-\sum_{a<k\le b}\frac1{k-t_0}=(b-a)\sum_{t=t_0}^{t_1}\frac1{(a-t)(b-t)}\quad(a,b\notin[t_0,t_1]),
\]
gives the entries modulo $\ell$: low--low $\bar x\,\delta_{pq}+A_{pq}$; low--middle $-\frac16\cdot\frac1{p-q}$; middle--middle
$-\frac12\delta_{pq}$; low--top and middle--top $0$. Here, in terms of $u_y=(1/(a+y))_{a<K}$,
\[
 A=\kappa\Bigl(hI+\sum_{n\in[K,T)}u_{\ell-n}u_{\ell-n}^{\top}\Bigr)+\tau_0\sum_{n\in[T,\ell)}u_{\ell-n}u_{\ell-n}^{\top},\qquad
 \kappa=-\tfrac1{18},\quad\tau_0=\tfrac1{36},
\]
and $h=\frac12\bigl(\sum_{y=1}^{3t}y^{-2}-2\sum_{y=3t+1}^{4t}y^{-2}\bigr)$ in $\F_\ell$ (note $\ell-K=4t$). The middle and top nodes enter with
weights of opposite sign. The identity for $A$ is the reflection $n=\ell-y$ of the two windows above, and the cross coefficient $\frac16$,
the middle diagonal $\frac12$ and $\kappa$ satisfy $(\frac16)^2=\frac12\cdot\frac1{18}$, which makes the middle level behave as in the
two-level case; the top level does not, and only enters through $A$. The Schur complement of
the middle block then gives, as in Proposition~\ref{prop:red}, bordered matrices
$\mathcal B(z,\;N_{\rm mid}+\frac{\tau_0}{\kappa}N_{\rm top})$, where $N_{\rm top}$ is the Gram matrix of the top columns, which are
the $w_j$ with $j\in[1,T-1]$, and $N_{\rm mid}$ that of the absent middle columns: $w_T$ for the first member and none for the
second. Hence the two members give
\[
\begin{gathered}
 \mathcal B\bigl(z,\ cN\bigr)\quad\text{and}\quad\mathcal B\bigl(z,\ cN+ww^{\top}\bigr),\\ N=\sum_{j=1}^{S-1}w_jw_j^{\top},\ w=w_S,\ S=T,\ c=\frac{\tau_0}{\kappa}=-\frac12 .
\end{gathered}
\]
Dividing the upper-left block by $c$ (and rescaling the border) turns this into the pair
$\mathcal B(z/c,N)$, $\mathcal B(z/c,N+c^{-1}ww^{\top})$, which is the situation of Lemma~\ref{lem:primal} with $\tau=c^{-1}=-2\ne0$.
The hypotheses \eqref{eq:nvrange} hold with $S=T$: $r+2m+S=\ell$, and $S=3t+1\ge K-m$. So Theorem~\ref{thm:good} and
Proposition~\ref{prop:delta} apply verbatim, all bases are at most $S+K+2m-1=\ell-r+K-1<\ell$, and we obtain the analogue of
Theorem~\ref{thm:nv}: if $\ell\nmid\den x_0$, one of the two members does not vanish at $x_0$. No information on the block
$M[X',Y']$ is needed.

\emph{Existence of pairs.} For $\ell=6t+1$ we need integers $r,m\ge1$ with $r+2m=3t$, $2m+1\le t$ and an admissible relation
between $r$ and $m$. With $r\le8m\le r+W$ this means $3t\le10m\le3t+W$, which is solvable for all $t$ if and only if $W\ge9$; for
$W=8$ it fails exactly when $\ell\equiv43\pmod{60}$. We therefore use $r\le8m\le r+9$ for $L(2,\chi_{-3})$; then every prime
$\ell\equiv1\pmod6$ with $\ell\ge61$ has a pair, with $5r\le2\ell\le5r+11$.

%% =====================================================================
% New section: bounded low-node deletions for the chi frame.
% This file deliberately leaves the preceding pair argument unchanged.
\section{Uniform nonvanishing after bounded low-node deletion}\label{sec:minor-nv}

For the $\chi$ frame, allowing a bounded number of common nodes to be omitted
removes the restriction $r+2m=3t$ in the preceding pair construction. The
borders remain fixed. We use the rational kernel $M$ of
Section~\ref{sec:construction}. For $J\subseteq[0,r)$ put
\[
 X'=\{r+2i:0\le i<m\},\qquad Y'=\{r+2i+1:0\le i<m\},
\]
and put $\xi_\chi=\tfrac14L(2,\chi_{-3})$. Let $\Phi_J(x)$ have rows $J\sqcup X'$, columns $J\sqcup Y'$, and entries
$M[a,b]+x\mathbf1[a=b\in J]$. Write $P_J(x)=\det\Phi_J(x)$.

\begin{theorem}\label{thm:minor-nv}
Let $\ell=6t+1$ be prime, and suppose
\begin{equation}\label{eq:minor-room}
 m\ge1,\qquad K:=2t+1\le r,\qquad r+2m\le3t.
\end{equation}
If $x\in\Q$ and $\ell^2\nmid\den x$, there is a set $I\subseteq[0,K)$,
$|I|\le3$, such that, for $J=[0,r)\setminus I$,
\[
 P_J(x)\ne0.
\]
\end{theorem}

The proof bounds the dimension of a kernel, rather than asserting simplicity
of its eigenvalues. Conditions~\eqref{eq:minor-room} imply $2m+1\le t$, so
$t\ge3$ and $\ell\ge19$. All occurrences of $2,3,6,18,36$ below are therefore
units modulo $\ell$.

\subsection{Reflected coordinates and the signed moment pencil}
Work first over $k=\F_\ell$. The physical node $b$ and its reflected coordinate
$y=\ell-b$ will be kept distinct. Put
\[
 S=3t+1,\qquad g=\ell-r-2m,\qquad d=K-m,\qquad h=r-K.
\]
Then $1<S\le g$, $3\le d<K$, $d\le g$, and
$K+g+2m+h=\ell$. For a reflected coordinate $y$ define
\[
 c_y=\bigl((a+y)^{-1}\bigr)_{0\le a<K},\qquad
 N=\sum_{y=1}^{g}\epsilon_y c_yc_y^{\top},\qquad
 \epsilon_y=\begin{cases}1&y<S,\\-2&y\ge S.\end{cases}
\]
Reversing the order of each physical border gives the reflected border
coordinates $g+2,g+4,\ldots,g+2m$ and
$g+1,g+3,\ldots,g+2m-1$. Denote their column matrices by $W_X,W_Y$, respectively,
and set
\[
 \Gamma=W_X^{\top}W_Y,\qquad
 B_z=\begin{pmatrix}N-zI_K&W_Y\\W_X^{\top}&0\end{pmatrix}.
\]
Proposition~\ref{prop:delta} gives $\det\Gamma\ne0$: its range condition is
$g+K+2m=\ell-r+K\le\ell$.

Introduce the polynomials
\begin{align*}
 D(y)&=\prod_{a=0}^{K-1}(y+a),&
 B_X(y)&=\prod_{i=0}^{m-1}(y-g-2i-2),\\
 B_Y(y)&=\prod_{i=0}^{m-1}(y-g-2i-1),& R(y)&=B_X(y)B_Y(y).
\end{align*}
For $0\le i<d$ define residue-basis matrices by
\begin{equation}\label{eq:minor-residues}
 (V_X)_{a i}=\frac{B_X(-a)(-a)^i}{D'(-a)},\qquad
 (V_Y)_{a i}=\frac{B_Y(-a)(-a)^i}{D'(-a)}.
\end{equation}
All displayed denominators and all $B_X(-a),B_Y(-a)$ are nonzero.
Partial fractions give, for $\deg p<d$,
\begin{equation}\label{eq:minor-pf}
 \sum_{a=0}^{K-1}\frac{B_X(-a)p(-a)}{D'(-a)(y+a)}
       =\frac{B_X(y)p(y)}{D(y)},
\end{equation}
and likewise for $Y$. Thus the columns of $V_X$ lie in
$\ker W_X^{\top}$. They are independent, since evaluation of a polynomial
of degree $<d\le K$ at the $K$ distinct points $-a$ is injective.
Since $\Gamma$ is invertible, $W_X,W_Y$ have rank $m$; consequently $V_X,V_Y$
are bases of $\ker W_X^{\top},\ker W_Y^{\top}$.

The compressed matrix is the symmetric Hankel pencil
\[
 A_z=V_Y^{\top}(N-zI_K)V_X=G-zH,
\]
whose entries, indexed by $0\le i,j<d$, are
\begin{align}
 H_{ij}&=\sum_{u\in\{0,-1,\ldots,1-K\}}
             \frac{R(u)}{D'(u)^2}u^{i+j},\label{eq:minor-H}\\
 G_{ij}&=\sum_{y=1}^{g}\epsilon_y\frac{R(y)}{D(y)^2}y^{i+j}.
 \label{eq:minor-G}
\end{align}
These identities follow from~\eqref{eq:minor-residues} and
\eqref{eq:minor-pf}. In particular the same polynomial coordinates describe
the left and right kernels; this fact will be needed for principal deletion.

\subsection{A uniform corank bound}
\begin{lemma}\label{lem:minor-corank}
For every $z\in k$, $\dim_k\ker A_z\le3$.
\end{lemma}
\begin{proof}
Identify a coefficient vector in $\ker A_z$ with a polynomial $p$, $\deg p<d$.
Equations~\eqref{eq:minor-H}--\eqref{eq:minor-G} say that its first $d$
weighted moments vanish. To use these conditions, complete the root sets by
\[
 E(y)=\prod_{j=1}^{g}(y-j),\qquad
 C(y)=\prod_{j=g+2m+1}^{\ell-K}(y-j).
\]
The last product is $1$ when $h=0$. The roots of $D,E,R,C$ partition $k$, so
\begin{equation}\label{eq:minor-fullfield}
 DERC=y^\ell-y,\qquad (y^\ell-y)'=-1.
\end{equation}
Let $T=DE$. On its simple roots define
\[
 w_z(v)=\begin{cases}
 -zR(v)/D'(v)^2&D(v)=0,\\
 \epsilon_v R(v)/D(v)^2&E(v)=0,
 \end{cases}
\]
and form the polynomial, linear in $p$,
\[
 B_p(y)=\sum_{T(v)=0}w_z(v)p(v)\frac{T(y)}{y-v},\qquad
 F_p(y)=D(y)^2C(y)B_p(y).
\]
The first $d$ moments are zero. Expanding
$B_p/T=\sum_v w_z(v)p(v)/(y-v)$ formally at infinity therefore gives
\[
 \deg B_p\le K+g-d-1=g+m-1,\qquad \deg F_p\le\ell+d-1.
\]
Equivalently these degree bounds follow by comparing the first $d$ leading
coefficients in the finite polynomial sum defining $B_p$.
Evaluating at a root of $T$, and differentiating
\eqref{eq:minor-fullfield}, gives the exact identities
\begin{equation}\label{eq:minor-interpolation}
 \begin{array}{ll}
 F_p(u)=F_p'(u)=0,\quad F_p''(u)=2zp(u)&(D(u)=0),\\
 F_p(j)=-\epsilon_jp(j)&(1\le j\le g),\\
 F_p(j)=0&(C(j)=0).
 \end{array}
\end{equation}
For example $B_p(u)=-zR(u)E(u)p(u)/D'(u)$ at a root $u$ of $D$;
then $F_p''(u)=2D'(u)^2C(u)B_p(u)=2zp(u)$ by
$D'(u)E(u)R(u)C(u)=-1$. The high-node identity follows in the same way.

Consider the finite-difference Wronskian
\[
 W_p(y)=F_p(y+1)p(y)-F_p(y)p(y+1).
\]
At each low adjacency $-K+1,\ldots,-1$, its value and first derivative
vanish. Its second derivative also vanishes, by the equality
$2zp(y+1)p(y)-2zp(y)p(y+1)=0$. Hence these are triple roots.
At the high adjacencies $1,\ldots,S-2$ and $S,\ldots,g-1$ its value
vanishes because the two weights agree. It vanishes at every root of $C$:
for the last such root, $\ell-K$, the next point is the first low root
$\ell-K+1$. Thus
\begin{equation}\label{eq:minor-fixedfactor}
 Q_0(y):=\prod_{a=1}^{K-1}(y+a)^3
       \prod_{j=1}^{S-2}(y-j)\prod_{j=S}^{g-1}(y-j)C(y)
       \quad\hbox{divides } W_p(y).
\end{equation}
The supports are disjoint, and
\[
 \begin{gathered}
 \deg Q_0=3(K-1)+(g-2)+h=\ell+2d-5,
 \\ Q_0(S-1)\ne0,\qquad W_p(S-1)=3p(S)p(S-1).
 \end{gathered}
\]

We also need that $W_p\ne0$ for nonzero $p$ in the kernel.
If $W_p=0$, reduce $F_p/p=A/B$ to coprime polynomials with $B$ monic.
Then $A(y+1)B(y)=A(y)B(y+1)$ implies $B\mid B(y+1)$, and hence
$B(y+1)=B(y)$. But $\deg B\le\deg p<\ell$; a nonconstant polynomial
of degree $e<\ell$ has a nonzero leading term of degree $e-1$ in its
first difference. Thus $B$ is constant. The polynomial quotient $F_p/p$
is translation invariant and has degree $<2\ell$. Its positive degree,
if any, must be $\ell$ by the same leading-coefficient argument.
Subtracting a multiple of $y^\ell-y$ therefore shows
\[
 F_p=p\bigl(a(y^\ell-y)+b\bigr).
\]
Some low node $u$ has $p(u)\ne0$, since $\deg p<K$. The first two
identities in~\eqref{eq:minor-interpolation} force $b=0$ and $a=0$.
Then the high-node identity forces $p$ to vanish at $g>\deg p$ points,
a contradiction.

Finally, if the kernel had dimension at least four, its three linear
conditions
\[
 [y^{d-1}]p=[y^{d-2}]p=0,\qquad p(S)=0
\]
would admit a nonzero solution. This solution has $\deg p\le d-3$.
The leading terms in the two products defining $W_p$ cancel, so
$\deg W_p\le\deg F_p+\deg p-1\le\ell+2d-5=\deg Q_0$.
By~\eqref{eq:minor-fixedfactor}, $W_p$ is a scalar multiple of $Q_0$;
its value at $S-1$ makes that scalar zero. This contradicts the preceding
paragraph and proves the lemma.
\end{proof}

\subsection{Selecting a nonzero principal minor}
\begin{lemma}\label{lem:minor-rescue}
For every $z\in k$, deleting at most three low principal coordinates from
$B_z$ leaves a nonsingular matrix.
\end{lemma}
\begin{proof}
Let $\mathcal P_z$ be the polynomial kernel of $A_z$, and put
$k_0=\dim\mathcal P_z\le3$. A right kernel vector of $B_z$ has low
part $v\in\ker W_X^{\top}$, hence $v=V_Xp$ uniquely. Its remaining
equation is $(N-zI)v\in\operatorname{im}W_Y$, which is equivalent to
$V_Y^{\top}(N-zI)V_Xp=0$. The border coordinates are then unique,
since $\Gamma$ is invertible. Conversely every $p\in\mathcal P_z$
lifts to such a kernel vector. The transpose has the corresponding
parametrization with low part $V_Yp$; the polynomial space is the same
because $A_z$ is symmetric.

Evaluation $p\mapsto(p(-a))_{a<K}$ is injective on $\mathcal P_z$.
Choose $k_0$ coordinates on which it is an isomorphism, and let $I$ be
the corresponding low indices. By~\eqref{eq:minor-residues}, restriction
of either full kernel to $I$ is an isomorphism: it differs from polynomial
evaluation only by nonzero diagonal factors.

Here is the elementary principal-minor criterion being used. Suppose a
square matrix has injective restriction of its right kernel to $I$ and
surjective restriction of its left kernel to $I$. A kernel vector of
the complementary principal matrix, extended by zero on $I$, is sent by
the full matrix to a vector supported on $I$. Pairing with left kernel
vectors shows that this image is zero. Right-kernel injectivity then
shows the extended vector is zero. Thus the complementary principal
matrix is nonsingular. Apply this criterion to $B_z$.
\end{proof}

\subsection{Return to the physical determinant}
Suppose first $\ell\nmid\den x$, and write $\bar x$ for its reduction.
Let
\[
 c_0=\sum_{y=1}^{3t}y^{-2}-2\sum_{y=3t+1}^{4t}y^{-2},\quad
 \kappa=-\frac1{18},\quad \tau_0=\frac1{36},\quad h_0=\frac{c_0}{2}.
\]
The three-level entry calculation of Section~\ref{ssec:nvother} applies
under~\eqref{eq:minor-room}; it does not require equality $r+2m=3t$.
The physical absent middle nodes are $[r+2m,S)$, which reflect to $[S,g]$;
the top nodes reflect to $[1,S)$. Their weighted Gram matrix is therefore
\[
 N_{\rm phys}=\sum_{y=S}^{g}c_yc_y^{\top}
              -\frac12\sum_{y=1}^{S-1}c_yc_y^{\top}=-\frac12 N.
\]
Consequently the physical reduced low block is exactly
\begin{equation}\label{eq:minor-physical}
 (h_0+\bar x/\kappa)I+N_{\rm phys}
       =-\frac12\bigl(N-zI\bigr),\qquad z=c_0-36\bar x.
\end{equation}
Both border orders are reversed simultaneously to obtain $W_X,W_Y$.
Scaling a low block by a nonzero scalar $c$ preserves bordered
nonsingularity, since
\[
 \begin{pmatrix}cI&0\\0&I\end{pmatrix}
 \begin{pmatrix}M_0&W_Y\\W_X^{\top}&0\end{pmatrix}
 \begin{pmatrix}I&0\\0&c^{-1}I\end{pmatrix}
 =\begin{pmatrix}cM_0&W_Y\\W_X^{\top}&0\end{pmatrix}.
\]
This also holds after any low principal restriction. Thus
Lemma~\ref{lem:minor-rescue} supplies a nonzero physical bordered minor.

For completeness, this minor transfers to the rational matrix without
assuming that the retained common set contains every low node.
Start with the original geometry and multiply its low rows and columns
by $1$, its middle rows and columns by $\ell$, and negate the matrix.
For any retained low set $U\subseteq[0,K)$ its reduced principal
restriction has the same Schur identity as the full matrix. Writing
$M_{\rm phys}$ for the left-hand side of~\eqref{eq:minor-physical},
\begin{equation}\label{eq:minor-schur-factor}
 \det Z_U=(1/2)^{r-K+m}(1/18)^{|U|}
 \det\begin{pmatrix}
 M_{\rm phys}[U,U]&W_Y[U,*]\\
 W_X[U,*]^{\top}&0
 \end{pmatrix}.
\end{equation}
Indeed the invertible middle common block is unchanged by low restriction;
the harmonic identity restricts entrywise, and the border row and column
operations are unchanged. Every prefactor is a unit. A zero rational
principal determinant would remain zero after the rational diagonal
scalings and after reduction, contradicting~\eqref{eq:minor-schur-factor}.
Taking $U=[0,K)\setminus I$ proves the theorem in this case.

It remains to treat $\ell\mid\den x$ but $\ell^2\nmid\den x$.
Coprimality of the reduced numerator and denominator shows that
$u=\overline{\ell x}$ is nonzero. Use the undeleted common set $J=[0,r)$.
Multiply \emph{every} row of $\Phi_J(x)$ by $\ell$ and additionally
multiply each middle common column by $\ell$. With rows and columns
ordered as low, middle common, border, the resulting integral matrix
reduces to
\[
 \begin{pmatrix}
 uI_K&0&-\tfrac16W_Y\\
 -\tfrac16W_C^{\top}&-\tfrac12I_{r-K}&0\\
 -\tfrac16W_X^{\top}&0&0
 \end{pmatrix}.
\]
Here $W_C$ consists of the middle common columns. In verifying the zero
blocks one uses integrality of distinct middle--middle entries, rather
than merely integrality after multiplication by $\ell^2$; it follows
directly from the off-diagonal divided-difference formula, whose interval
contains neither singular index $K$ nor $S$.
Reorder the middle common block last. Its determinant is $(-1/2)^{r-K}$,
and the remaining bordered determinant is
\[
 \det\begin{pmatrix}uI_K&-\tfrac16W_Y\\-\tfrac16W_X^{\top}&0\end{pmatrix}
       =u^{K-m}(-1/36)^m\det\Gamma\ne0.
\]
Thus $P_{[0,r)}(x)\ne0$. This proves
Theorem~\ref{thm:minor-nv} in the remaining case.

\subsection{Content, decay, and the original degree budget}
The enlarged family consists of all $P_J$ with $|[0,r)\setminus J|\le3$.
For admissible $r\le8m\le r+9$ it retains the constants
\[
 \Lambda=\frac{10339}{3072},\qquad D=3.7251,\qquad c=6.6510.
\]
To see the content assertion, use the entrywise tropical denominator
bound from Section~\ref{sec:content}, before any determinant cancellation.
Any assignment for a retained principal submatrix extends to one for
the original matrix by fixing the omitted common coordinates. All
weights are nonnegative, so the smaller assignment value cannot be
larger. Write $d_{r,m}$ for this common integer denominator. It clears every $P_J$, and
its logarithm is at most $(\Lambda+o(1))r^2$, uniformly in $J$.

The energy estimates of Section~\ref{sec:decay} already apply to every
subset of common nodes, with the borders retained. If $k\le3$ common
nodes are forced to be absent, expansion at $\xi_\chi+v$ has terms
\[
 v^{r-k-|A|}\det\bigl(\omega\Lw_{F_\chi}[A\sqcup X',A\sqcup Y']\bigr),
 \qquad A\subseteq J.
\]
For $|v|\le e^{-cr}$, use
$|v|^{r-k-|A|}\le e^{-cr(r-k-|A|)}$. Replacing the original exponent
$r-|A|$ by $r-k-|A|$ therefore costs at most $e^{crk}$ in the termwise
upper bound, also at $v=0$. This is
a logarithmic loss of at most $3cr=O(r)$, and the number of terms is
still at most $2^r$. It is absorbed by the existing $o(r^2)$ errors.
Hence $|P_J(\xi_\chi+v)|\le\exp(-(D-o(1))r^2)$ on the same window,
uniformly in $J$; no comparison of possibly cancelling determinants is
needed.

Finally $\deg P_J\le|J|\le r$. We keep the declared degree budget equal
to the original $r$ for every member, so $q^r$ times the common clearing
denominator times $P_J(p/q)$ is an integer. Neither the degree scale nor
the exponential window changes when the nonzero member is selected.


\section{Analytic decay}\label{sec:decay}

\begin{theorem}[decay and window]\label{thm:decay}
For admissible geometries of scale $r$ and real $t$ with $|t|\le e^{-cr}$,
\[
\begin{gathered}
 |P_{\mathcal G}(\xi+t)|\le\exp\bigl(-(D-o(1))r^2\bigr),\\
 (D,c)=\begin{cases}(3.9028,\ 7.6444)&\text{for }\lambda,\lambda',\\ (3.7251,\ 6.6510)&\text{for }L(2,\chi_{-3}).\end{cases}
\end{gathered}
\]
\end{theorem}

We describe the proof for $\lambda$; the other two cases differ only in the reflection formula and the certificate
data.

\subsection{A contour representation}
Hermite's formula gives $\Lw_F[a,b]=\frac1{2\pi i}\oint F(z)v_a(z)v_b(z)\,dz$ with $v_a(z)=1/(z-a)$, the contour enclosing
the nodes and no pole of $F$. The reflection formulas $\psi(1+z)=\psi(-z)-\pi\cot\pi z$ and
$\psi(\frac34+z)=\psi(\frac14-z)+\pi\cot(\frac\pi4-\pi z)$ split $F=\tilde F+E$, where $\tilde F$ is analytic in the left
half-plane and $O(1/|z|)$ there, and
\[
 E(z)=-\pi\bigl[\cot\pi z+\cot(\tfrac\pi4-\pi z)\bigr]
\]
is $1$-periodic with $|E(x+iy)|\le C_Ee^{-2\pi|y|}$ for $|y|\ge1$. Let $N_0=\lfloor\frac7{50}r\rfloor$, $x^*=-N_0-\frac38$, and let
$\Gamma$ be the two rays $x^*+u(-1\pm3i)$, $u\ge0$. Moving the contour to $\Gamma$ picks up the residues of $F$ at the
$2N_0$ poles $-(j+\frac34)$ (residue $+1$) and $-(j+1)$ (residue $-1$), $j<N_0$; the part of the integrand with $\tilde F$
integrates to $0$ over $\Gamma$. Hence, for all integers $a,b\ge0$,
\begin{equation}\label{eq:gram}
 \Lw_F[a,b]=\int v_av_b\,dm,\qquad m=-\sum_{\mu}\epsilon_\mu\,\delta_{-\mu}+\frac1{2\pi i}E(z)\,dz\big|_\Gamma ,
\end{equation}
a Gram representation with respect to a complex measure $m$ with finitely many atoms and an exponentially small
density on $\Gamma$.

\subsection{Two-family Andr\'eief bound}
For a set $S\subseteq C$ of common nodes let $R_S=S\sqcup X'$, $Q_S=S\sqcup Y'$, $n=|S|+m$, and
$\mathcal M_S=\omega\,\Lw_F[R_S,Q_S]$. By multilinearity in the diagonal,
\begin{equation}\label{eq:expand}
 P_{\mathcal G}(\xi+t)=\sum_{S\subseteq C}t^{|C|-|S|}\det\mathcal M_S .
\end{equation}
Andr\'eief's identity \cite{Forrester2018Andreief} for the two families $(v_a)_{a\in R_S}$ and $(v_b)_{b\in Q_S}$ gives
\[
 \det\Lw_F[R_S,Q_S]=\frac1{n!}\int\det[v_a(z_i)]\,\det[v_b(z_i)]\prod_{i=1}^ndm(z_i),
\]
and Cauchy's determinant evaluates both factors. Taking absolute values (see \cite{SaffTotik} for the
logarithmic potential theory behind what follows),
\begin{multline*}
 |\det\Lw_F[R_S,Q_S]|\le\frac1{n!}\int\prod_{i<i'}|z_i-z_{i'}|^2\,\\
 \prod_{\mathcal F\in\{R_S,Q_S\}}\frac{\prod_{a<a'\in\mathcal F}|a'-a|}{\prod_{i,\,a\in\mathcal F}|z_i-a|}\ \prod_id|m|(z_i).
\end{multline*}
The integrand is the exponential of minus a discrete logarithmic energy: the free points $z_i$ repel each other with
weight $1$, each node family repels itself with weight $\frac12$ and attracts the free points with weight $1$, and the
density $|E|\approx e^{-2\pi|\Im z|}$ acts as an external field. Because the rows and columns differ only in the
interleaved blocks, the two node families have the same macroscopic density: after scaling by $r$, both are
\[
 \nu_\alpha=\mathrm{Leb}|_{[0,1]}+\tfrac12\mathrm{Leb}|_{[1,1+2\alpha]},\qquad\alpha=\tfrac18,
\]
of mass $1+\alpha$. Its logarithmic energy has the closed form
\begin{multline*}
 I(\nu_{1/8})=\iint-\log|u-u'|\,d\nu_{1/8}(u)\,d\nu_{1/8}(u')\\
 =\frac{243}{128}-\frac{\log2}{32}-\frac{25}{32}\log\frac54=1.70244\ldots
\end{multline*}

\subsection{The energy certificate}
Fix $\rho=1/1024$ and the regularised kernel $k_\rho(z,w)=-\frac12\log(|z-w|^2+\rho^2)\le-\log|z-w|$.

\begin{lemma}[conditional positive definiteness]\label{lem:cpd}
For points $p_1,\dots,p_N\in\C$ and real weights $q_\alpha$ with $\sum q_\alpha=0$,
$\sum_{\alpha,\beta}q_\alpha q_\beta k_\rho(p_\alpha,p_\beta)\ge0$.
\end{lemma}
\begin{proof}
Frullani's integral $\log a=\int_0^\infty(e^{-t}-e^{-at})\,dt/t$ and $\sum q_\alpha=0$ give
$\sum q_\alpha q_\beta(-\log(|p_\alpha-p_\beta|^2+\rho^2))=\int_0^\infty\frac{e^{-t\rho^2}}t\sum q_\alpha q_\beta e^{-t|p_\alpha-p_\beta|^2}dt$,
and Gaussian kernels are positive definite: expanding $e^{2t\langle p,p'\rangle}$ in powers and the binomial theorem
write the inner sum as a convergent sum of squares.
\end{proof}

Apply Lemma~\ref{lem:cpd} to the free points and removed nodes (weight $1/r$ each) against a \emph{certificate
measure} $\sigma=\sum_km_k\delta_{c_k}$ of total mass $1+\alpha$. This bounds the free--free energy from below by a
sum of one-point terms $U^\sigma_\rho(s)=\sum_km_kk_\rho(s,c_k)$ and the constant $I_\rho(\sigma)=\sum_{k,l}m_km_lk_\rho(c_k,c_l)$.
Riemann-sum comparisons replace the node sums by the potential $U^{\nu_\alpha}$ and the node self-energy by
$I(\nu_\alpha)$, up to $O(r\log r)$. With the external field one arrives at the function
\[
 g(s)=(1+\alpha)U^\sigma_\rho(s)-U^{\nu_\alpha}(s)+(1-\kappa_0)\pi|\Im s|,\qquad\kappa_0=10^{-3},
\]
on the scaled support of $|m|$ (the rays $\Gamma/r$, the segment $[-\frac7{50},0]$ containing the atoms, and, for removed
nodes, the segment $[0,1+2\alpha]$). The certificate consists of three finite families of inequalities, checked in
interval arithmetic:
\begin{enumerate}[label=(\alph*),leftmargin=*]
\item $g\ge m'$ on a cover of the rays by boxes, with an analytic tail bound beyond $u=\frac32$;
\item $g\ge\gamma_\ell$ on a partition of $[-\frac7{50},0]$ into cells; put $S_{\rm cert}=\allowbreak2\sum_\ell|{\rm cell}_\ell|\allowbreak\min(\gamma_\ell-m',0)$;
\item $g\ge g_{\min}$ on a cover of $[0,1+2\alpha]$.
\end{enumerate}
On a box, $U^\sigma_\rho$ is bounded below by moving each $c_k$ to the farthest corner, and $U^{\nu_\alpha}$ is bounded
above from its closed form and monotonicity. Then
\[
 D=2\bigl((1+\alpha)m'+S_{\rm cert}\bigr)-(1+\alpha)^2I_\rho(\sigma)+I(\nu_\alpha),\qquad c=2(m'-g_{\min}).
\]

\begin{proposition}[certified constants]\label{prop:cert}
For $\lambda$ and $\lambda'$ there is a measure $\sigma$ with $480$ atoms with dyadic coordinates and masses for which
\[
\begin{gathered}
 m'\ge2.2147,\quad S_{\rm cert}\ge-0.12218,\quad (1+\alpha)^2I_\rho(\sigma)\le2.53835,\\ g_{\min}\ge-1.6075,\quad
 I(\nu_{1/8})\ge1.7024457,
\end{gathered}
\]
whence $D\ge3.9028$ and $c\le7.6444$. For $L(2,\chi_{-3})$ the certificate of \S\ref{ssec:chicert} gives $D\ge3.7251$
and $c\le6.6510$.
\end{proposition}

All the inequalities are verified by the Lean kernel (\S\ref{sec:lean}); the data were generated with
interval arithmetic in Arb \cite{Arb} and then checked independently in Lean with rational arithmetic. For $\lambda$
the measure is assembled from a square-frame certificate and an extra family of atoms that represents the additional
density $\frac12\mathrm{Leb}|_{[1,1+2\alpha]}$; the function $g$ splits as the square-frame function plus a correction
term, which is certified on a coarse cover.

\subsection{The certificate for $L(2,\chi_{-3})$}\label{ssec:chicert}
For $F_\chi$ the reflection formula $\psi(z+a)=\psi(1-a-z)-\pi\cot\pi(z+a)$ gives
$F_\chi=\tilde F_\chi+E_\chi$ with
\[
 E_\chi(z)=\pi\Bigl[\tfrac13\cot\pi(z+\tfrac13)-\tfrac12\cot\pi(z+\tfrac12)+\tfrac16\cot\pi(z+1)\Bigr],
\]
which is $1$-periodic and, because the weights $\frac13-\frac12+\frac16$ sum to zero, decays like $e^{-2\pi|\Im z|}$. The
contour argument and the Andr\'eief bound are as above, with poles at $-(j+\frac13)$, $-(j+\frac12)$, $-(j+1)$. The square-frame
certificate for $F_\chi$ (a measure $\sigma_\chi$ of $699$ atoms, $\rho=1/2048$) gives lower bounds for the square-frame
function $g_\chi=U^{\sigma_\chi}_\rho-U^{\tau}+(1-\kappa_0)\pi|\Im s|$ on rays, segment and line. The hybrid certificate is
\[
 \sigma_{\rm h}=\frac{\sigma_\chi+\lambda}{9/8},\qquad
 \lambda=\sum_{k=0}^{31}\frac1{256}\,\delta_{t_k},\qquad t_k=1+\frac k{128}-\frac1{2048},
\]
so that $\lambda$ has mass $\frac18$ and represents the extra density $\frac12\mathrm{Leb}|_{[1,5/4]}$; the atom $t_k$ sits at the
left end $a_k=1+k/128$ of the $k$-th cell, shifted left by $\rho$. The function of the hybrid frame splits as
$g_{\rm h}=g_\chi+h$ with $h=U^{\lambda}_\rho-U^{\rm ex}$, where $U^{\rm ex}$ is the potential of $\frac12\mathrm{Leb}|_{[1,5/4]}$.
For $u$ in the $k$-th cell and $\Re s+\rho\le a_k$ one has
$(\Re s-t_k)^2+\rho^2\le(\Re s-u)^2$, hence $k_\rho(s,t_k)\ge-\log|s-u|$ on the cell, and integrating over the cells gives
$h(s)\ge0$ whenever $\Re s\le1-\rho$. On the short real interval $[1-\rho,1]$ only the first cell is affected, and an elementary
bound for $\int_0^\epsilon-\log v\,dv$ gives $h\ge0$ there too. Therefore every lower bound of the square-frame certificate holds for
$g_{\rm h}$, and only the energy $I_\rho(\sigma_{\rm h})$ is new: it consists of the energy of $\sigma_\chi$, the cross terms
($699\times32$ logarithms) and the self-energy of $\lambda$ ($32^2$ logarithms), summed exactly in rational interval
arithmetic. The result is $D\ge3.7251$ (the computed value is $3.725194$) and $c\le6.6510$; with $64$ extra atoms instead of
$32$ one would gain only $0.0008$ in $D$.

Since the admissible relation for $L(2,\chi_{-3})$ is $r\le8m\le r+9$, the decay estimate is proved uniformly for
$\alpha r\le m\le\alpha r+\frac98$, $\alpha=\frac18$: the comparison of node sums with $U^{\nu_\alpha}$ and $I(\nu_\alpha)$ is carried out
with a window parameter for the excess of $m$ over $\alpha r$, which only changes the $O(r\log r)$ error terms.

\subsection{Node removal and the window}
The argument applies to every subset $S\subseteq C$: removing $j=\beta r$ common nodes turns them into points of the
configuration that are not integrated over, and only common nodes can be removed, because a node removed from one
family only would leave an uncancelled term of weight $\frac12$. One obtains
\[
 |\det\mathcal M_S|\le\exp\bigl(-r^2(D-\beta c-o(1))\bigr)\qquad(\beta=|C\setminus S|/r),
\]
uniformly in $S$. Inserting this into \eqref{eq:expand} with $|t|e^{cr}\le1$ gives Theorem~\ref{thm:decay}: the sum has at
most $2^{|C|}=e^{O(r)}$ terms, each at most $e^{-(D-o(1))r^2}$. The second member is handled by viewing it as the pattern
$(r,m+1)$ with one column node moved by one place, at an energy cost $O(r)$.

%% =====================================================================
\section{Proof of Theorem~\ref{thm:main}}\label{sec:proof}

\subsection{$G\pm5\pi^2/48$}
Let $\xi=\lambda$ (the case $\lambda'$ is identical, with primes $\ell\equiv1\pmod 4$). Let
$\eta=D-\Lambda\ge3.9028-3.68274=0.22006$, let $\nu>c/\eta$, and suppose $|\xi-p/q|<q^{-\nu}$ with $q$ large; put
$L=\log q$. Choose $\delta>0$ with $\nu>c(1+\delta)/\eta$, let $r=\lceil(1+\delta)L/\eta\rceil$ and choose $m$ with
$r\le8m\le r+8$.

\emph{Choice of $\ell$.} By the prime number theorem in progressions, the primes $\ell\equiv3\pmod4$ in
$[\frac43(r+m)+1,4r-1]$ have product $\exp((\frac54+o(1))r)$ (as $m\le\frac r8+1$), which exceeds $q=e^L$ since $r\gg L$. So
one of them does not divide $q$. By Theorem~\ref{thm:nv} with $x_0=p/q$, there is $e\in\{0,1\}$ with
$P_{\mathcal G_e}(p/q)\ne0$, where $\mathcal G_e=\mathcal G(r,m,e)$.

\emph{Lower bound.} By Lemma~\ref{lem:integrality}, $|d_{\mathcal G_e}P_{\mathcal G_e}(p/q)|\ge q^{-(r+1)}=e^{-(r+1)L}$.

\emph{Upper bound.} Write $p/q=\xi+t$ with $|t|<e^{-\nu L}\le e^{-cr}$ (as $cr\le c(1+\delta)L/\eta+c<\nu L$ for large $q$).
By Theorems~\ref{thm:content} and~\ref{thm:decay},
$|d_{\mathcal G_e}P_{\mathcal G_e}(p/q)|\le e^{(\Lambda+o(1))r^2}e^{-(D-o(1))r^2}=e^{-(\eta-o(1))r^2}$.

\emph{Contradiction.} Comparing, $(\eta-o(1))r^2\le(r+1)L$; but $\eta r\ge(1+\delta)L$. Hence $\mu(\lambda)\le c/\eta$, and
\[
 \frac c\eta\le\frac{7.6444}{3.9028-3.68274}=34.737798\ldots<34.74 .
\]
Since $G-5\pi^2/48=2\lambda$ and $G+5\pi^2/48=2\lambda'$, this gives the first two bounds of Theorem~\ref{thm:main}.

\subsection{$L(2,\chi_{-3})$}
\emph{The earlier pair-family bound.}
Now the geometry is tied to the prime: for each prime $\ell\equiv1\pmod6$, $\ell\ge61$, there is an admissible pair of
members of scale $r$ with $5r\le2\ell\le5r+11$, so $r/\ell\to g=\frac25$, and Theorem~\ref{thm:nv} (in the form of
\S\ref{ssec:nvother}) applies when $\ell\nmid q$. Let $\eta=3.7251-\frac{10339}{3072}=0.359540\ldots$, $\nu>c(1/\eta+\frac45)$,
and choose $\delta$ small. The primes $\ell\equiv1\pmod 6$ in an interval $(Y,Y+(2+\delta)L]$ with $Y=\frac52(1+\delta)L/\eta$
have product $e^{(1+\delta/2+o(1))L}>q$, so one of them, $\ell$, does not divide $q$. The corresponding scale satisfies
\[
 (1+\delta)L/\eta-O(1)\le r\le\tfrac25\bigl(Y+(2+\delta)L\bigr)+O(1)=\Bigl(\frac{1+\delta}\eta+\frac{2(2+\delta)}5\Bigr)L+O(1),
\]
so $cr<\nu L$ for $\delta$ small, and the comparison above, with $|C|\le r+1$, gives a contradiction. Hence
\[
 \mu\bigl(\tfrac14L(2,\chi_{-3})\bigr)\le c\Bigl(\frac1\eta+gs\Bigr),\qquad g=\tfrac25,\ s=2,
\]
and $6.6510\,(1/0.3595401+0.8)=23.8194\ldots<23.82$. The valuation-one argument of Section~\ref{sec:minor-nv} also allows $s=1$, giving the retained bound $21.16$.

\emph{The bounded-deletion family.} Put
\[
 \eta=\frac{37251}{10000}-\frac{10339}{3072}=\frac{690317}{1920000},\qquad c=\frac{6651}{1000}.
\]
Let $\nu>c/\eta$, choose $\delta>0$ so that $c(1+\delta)/\eta<\nu$, put $L=\log q$ and
$r=\lceil(1+\delta)L/\eta\rceil$, and set $m=\lceil r/8\rceil$.
For all sufficiently large $r$, every prime $\ell\equiv1\pmod6$ in
\[
 \frac{51}{20}r\le\ell\le\frac{59}{20}r
\]
satisfies $K=(\ell+2)/3\le r$ and $r+2m\le(\ell-1)/2$.
Indeed $r+2m\le5r/4+7/4$ and the lower endpoint leaves a positive linear margin;
the upper endpoint gives $K\le59r/60+2/3\le r$ for large $r$.
The prime number theorem in this progression gives
\[
 \log\prod_{\substack{51r/20\le\ell\le59r/20\\\ell\equiv1\ (6)}}\ell^2
   =\left(\frac25+o(1)\right)r>L,
\]
since $2/5>\eta$. Thus some such prime satisfies $\ell^2\nmid q$.
By Theorem~\ref{thm:minor-nv}, some retained common set $J$ with at most three omissions has
$P_J(p/q)\ne0$. The original tropical denominator $d_{r,m}$ clears its coefficients, its degree is at most $r$,
and the estimates in Section~\ref{sec:minor-nv} give
\[
 q^{-r}\le |d_{r,m}P_J(p/q)|
 \le \exp\bigl(- (\eta-o(1))r^2\bigr)
\]
whenever $|\xi_\chi-p/q|<q^{-\nu}$; the required window holds because $cr<\nu L$.
This contradicts $\eta r\ge(1+\delta)L$. Consequently
\[
 \mu(\xi_\chi)\le\frac c\eta=\frac{12769920}{690317}
 =18.498631788\ldots<\frac{37}{2}.
\]
Rational scaling by four preserves the exponent, completing the proof of Theorem~\ref{thm:main}. \qed

\begin{remark}[relation to linear independence criteria]
The deduction is a one-dimensional version of the familiar scheme behind the criteria of Nesterenko
\cite{Nesterenko85} and their use by Ball and Rivoal \cite{BallRivoal01}: the nonzero integer
$q^{|C|}d_{\mathcal G}P_{\mathcal G}(p/q)$ plays the role of a small linear form with integer coefficients. The difference
is that the lower bound for the form comes from an arithmetic nonvanishing statement (Theorem~\ref{thm:nv}) rather than
from an analytic asymptotic, and that the size $r$ of the determinant is chosen in terms of $q$.
\end{remark}

\begin{remark}[where the constants enter]
For $\lambda$, $\partial\mu/\partial D=-c/\eta^2\approx-158$, so the third decimal of $D$ matters. The rank deficits of
\S\ref{sec:content} first lower the exponent from $50.0$ to $38.35$, and Section~\ref{sec:uniform-deficit} lowers it to $34.74$. The free choice of $\ell$ (Theorem~\ref{thm:nv} for all
$S$ in a range) removes a term $c\,g\,s$ which would otherwise be about $5.1$.
For $L(2,\chi_{-3})$, valuation-one nonvanishing first reduces the pair-family bound from $23.82$ to $21.16$;
bounded low-node deletion then removes the remaining prime-location cost and gives $18.50$.
With the present constants the limiting ratio is $18.498631788\ldots$. Reducing the fixed number of deletions,
or proving nonvanishing for a smaller fixed family, would not change this ratio.
\end{remark}

%% =====================================================================
\section{Formalization}\label{sec:lean}
\raggedbottom

All results of this paper are formalized in Lean~4 \cite{Lean4} with Mathlib \cite{Mathlib}. The development is
available at
\begin{center}
\FormalizationLink
\end{center}
and combines analytic and arithmetic proofs with generated rational data and kernel checks for the dual
certificates of Section~\ref{sec:content} and the decay certificates of Section~\ref{sec:decay}; numerical
computations alone are not used as proof. The constants are written out as series, so the main statements use no
definitions of the development: $G$ is $\sum_j(-1)^j/(2j+1)^2$ and $L(2,\chi_{-3})$ is
$\sum_k\bigl((3k+1)^{-2}-(3k+2)^{-2}\bigr)$. The main statements are stated in \lean{ChallengeHybrid.lean} and proved in
\lean{SolutionHybrid.lean} (namespace \lean{CatalanHybridChallenge}):
\begin{itemize}
\item \lean{lam_hybrid_measure_3474}: for every $e>0$, all sufficiently large $q\in\N$ and all $p\in\Z$ coprime to $q$,
 $q^{-(34.74+e)}<|G-5\pi^2/48-p/q|$ (Theorem~\ref{thm:main}); it is \lean{lam_hyb_measure_uniform};
\item \lean{mir_hybrid_measure_3474}: the same for $G+5\pi^2/48$; it is \lean{mir_hyb_measure_uniform};
\item \lean{chi3_hybrid_measure_1850}: the same for $L(2,\chi_{-3})$ with exponent $18.50$; it is
 \lean{chi3_hybrid_measure_minor};
\item \lean{p6_minorNV_square}, \lean{bandPencil_corank} and \lean{physical_low_rescue}:
 the bounded-deletion nonvanishing argument of Section~\ref{sec:minor-nv};
\item \lean{nv_hyb_nonvanishing_gen}: Theorem~\ref{thm:nv} and its analogues in \S\ref{ssec:nvother};
\item \lean{hyb_primal_false} and \lean{primal_pair}: the end of the nonvanishing chain and Lemma~\ref{lem:primal};
\item \lean{lam_hybrid_decay}, \lean{p6_hybrid_decay}: Theorem~\ref{thm:decay} for $\lambda$ and $L(2,\chi_{-3})$;
\item \lean{profile_both_eps}, \lean{p6_content}: Theorem~\ref{thm:content} for $\lambda$ and $L(2,\chi_{-3})$.
\end{itemize}
The solution has no \lean{sorry} and no \lean{native_decide}, and the command \texttt{\#print axioms} reports only the
three standard axioms \lean{propext}, \lean{Classical.choice} and \lean{Quot.sound}. The file
\lean{comparator-hybrid.json} is a configuration for \lean{lake comparator} stating that \lean{SolutionHybrid} proves the statements of
\lean{ChallengeHybrid} with these axioms only. We have not run Comparator, which requires Linux; the axiom check above was done with
\texttt{\#print axioms}. The prime number theorem in arithmetic progressions (moduli $4$ and $6$)
enters the content asymptotics and the selection of the prime $\ell$. The formal source for it is the PNT+ project
\cite{PNTPlus}, through OpenAI's workspace and compatibility additions, whose Wiener--Ikehara theorem we use.

\begingroup
\noindent The following file groups implement the indicated parts of the proof.\par\medskip
\noindent \lean{CatalanCore/Rect}\par
Geometries, the matrices $\Phi_{\mathcal G}$, Lemma~\ref{lem:integrality}, pair families and admissible relations.\par\medskip
\noindent \lean{CatalanCore/Arith/Hybrid}\par
Gauss valuations, Lemma~\ref{lem:trop}, the level majorant, Proposition~\ref{prop:dual}, prime sums over linear pieces,
the rank-deficit interface (Lemma~\ref{lem:rank}) and bordering.\par\medskip
\noindent \lean{CatalanFrames/Lambda/Hyb*}, \lean{CatalanFrames/LambdaMirrorHybrid}, \lean{CatalanFrames/P6Hybrid}\par
Theorem~\ref{thm:content} and Theorem~\ref{thm:deficits} for the three constants.\par\medskip
\noindent \lean{CatalanFrames/LambdaHybrid/NV}\par
Proposition~\ref{prop:red} and its three-level analogue, Lemma~\ref{lem:primal} (\lean{Primal}), Theorem~\ref{thm:nv}
(\lean{FinalNV}).\par\medskip
\noindent \lean{CatalanCore/Star/Leonard}\par
Proposition~\ref{prop:delta}, Proposition~\ref{prop:bispectral}, (K1)--(K5), the diagonal and Darboux steps, the
positivity argument, Lemma~\ref{lem:bridge}, Theorem~\ref{thm:good} and the specialisation $\theta\to0$.\par\medskip
\noindent \lean{CatalanCore/Decay/Rect}, \lean{CatalanCore/Cert/SigmaHyb}, \lean{CatalanCore/Cert/SigmaP6H}\par
The contour identity, the two-family Andr\'eief bound, Lemma~\ref{lem:cpd}, node removal, and the certificates of
Proposition~\ref{prop:cert} and \S\ref{ssec:chicert}.\par\medskip
\noindent \lean{CatalanCore/Measure}, \lean{CatalanFrames/*/Assembly*}\par
The deduction of Section~\ref{sec:proof}.\par\medskip
\endgroup

The proofs were checked with Lean \texttt{v4.35.0-rc2} and the matching Mathlib release. The vendored files from OpenAI's repository and from the PNT+ project were converted to Lean's module
system and ported to the current toolchain; the repository's README and provenance documents describe the changes.

\section*{Disclosure of AI use}

The original development was produced by the AI systems Claude Opus~5.5 and Claude Sonnet~5.5
(Anthropic): the mathematics, including the rectangular frames, the content bounds and rank deficits, the
nonvanishing argument and the decay certificates, and all proofs; the Lean formalization; the numerical
computations and the certificate generators; the literature search; and the text of this paper. Claude worked in a
series of conversations at the request of the author, who posed the problem, directed the work, and is responsible
for this submission. The repository infrastructure, vendored formalizations and certificate tooling of the
companion project \cite{Litsey2026}, which this work reuses, were developed with OpenAI Codex as described there.
The subsequent improvement for $L(2,\chi_{-3})$ from $23.82$ to $18.50$, including valuation-one reduction,
the finite-field corank argument, bounded-deletion transfer, their Lean formalization and the corresponding
additions to this paper, was developed with OpenAI Codex. The earlier exposition is retained except where the
stronger theorem requires changes. The subsequent uniform rescaled-core argument for both
$G\pm5\pi^2/48$, the finite interval profile and its Lean formalization, and the improvement from $38.35$
to $34.74$ were also developed with OpenAI Codex.
The underlying determinant method and the irrationality proof for $G$ are OpenAI's \cite{OpenAI2026}. We use Lean~4
\cite{Lean4}, Mathlib \cite{Mathlib}, the PNT+ project \cite{PNTPlus}, and Arb \cite{Arb} (through
\texttt{python-flint}) for interval arithmetic in generating certificate data, which Lean re-checks. The formal
verification in Lean is intended as an independent check of every step of the argument above.

% Bibliography generated by BibTeX (amsplain) from refs.bib; inlined for upload.
\providecommand{\bysame}{\leavevmode\hbox to3em{\hrulefill}\thinspace}
\providecommand{\MR}{\relax\ifhmode\unskip\space\fi MR }
% \MRhref is called by the amsart/book/proc definition of \MR.
\providecommand{\MRhref}[2]{%
  \href{http://www.ams.org/mathscinet-getitem?mr=#1}{#2}
}
\providecommand{\href}[2]{#2}
\begin{thebibliography}{10}

\bibitem{AAR99}
George~E. Andrews, Richard Askey, and Ranjan Roy, \emph{Special functions},
  Encyclopedia of Mathematics and its Applications, vol.~71, Cambridge
  University Press, Cambridge, 1999.

\bibitem{BallRivoal01}
Keith Ball and Tanguy Rivoal, \emph{Irrationalit\'e d'une infinit\'e de valeurs
  de la fonction z\^eta aux entiers impairs}, Invent. Math. \textbf{146}
  (2001), no.~1, 193--207.

\bibitem{Beukers79}
Frits Beukers, \emph{A note on the irrationality of $\zeta(2)$ and $\zeta(3)$},
  Bull. London Math. Soc. \textbf{11} (1979), no.~3, 268--272.

\bibitem{Brown26}
Francis Brown, \emph{Mellin transforms, transfinite diameter and rational
  approximations of integrals}, 2026, Preprint, 22 April 2026.

\bibitem{CDT2024}
Frank Calegari, Vesselin Dimitrov, and Yunqing Tang, \emph{The linear
  independence of $1$, $\zeta(2)$, and $l(2,\chi_{-3})$}, 2024.

\bibitem{CDT2025}
\bysame, \emph{Arithmetic holonomy bounds and effective {D}iophantine
  approximation}, 2025, Theorem 5.1 gives $\mu(L(2,\chi_{-3}))\le 24781$
  (checked at the source by S.~Litsey).

\bibitem{Carreira2026}
Marcos Costa Santos~Carreira, \emph{General designs and an irrationality
  exponent for {C}atalan's constant}, Preprint v1, October 12, 2026,
  \href{https://doi.org/10.5281/zenodo.23326325}{doi:10.5281/zenodo.23326325};
  sources at \url{https://github.com/MarcosCarreira/catalan-designs}, 2026.

\bibitem{Davenport}
Harold Davenport, \emph{Multiplicative number theory}, third ed., Graduate
  Texts in Mathematics, vol.~74, Springer, New York, 2000, Revised and with a
  preface by H.~L.~Montgomery.

\bibitem{Lean4}
Leonardo de~Moura and Sebastian Ullrich, \emph{The {L}ean 4 theorem prover and
  programming language}, Automated Deduction -- CADE 28, Lecture Notes in
  Computer Science, vol. 12699, Springer, 2021, pp.~625--635.

\bibitem{Eskandari2026}
Payman Eskandari, \emph{Improved rational approximations to {C}atalan's
  constant}, 2026, Preprint, 22 September 2026.

\bibitem{EMN2025}
Payman Eskandari, Kumar Murty, and Yusuke Nemoto, \emph{Mixed motives and
  linear forms in the {C}atalan constant}, 2025, Preprint, October 2025 (v3: 29
  August 2026).

\bibitem{Fauzan2026}
Aabir Fauzan, \emph{$\zeta(5)$ is irrational}, Preprint v1, September 17, 2026,
  Zenodo,
  \href{https://doi.org/10.5281/zenodo.22826419}{doi:10.5281/zenodo.22826419},
  2026.

\bibitem{Forrester2018Andreief}
Peter~J. Forrester, \emph{Meet {A}ndr\'eief, {B}ordeaux 1886, and {A}ndreev,
  {K}harkov 1882--83}, 2018.

\bibitem{Giraudin26}
David~Niedbala Giraudin, \emph{A new upper bound for the irrationality exponent
  of $\zeta(2)$}, 2026, Preprint, 2 October 2026.

\bibitem{GruenbaumHaine97}
F.~Alberto Gr\"unbaum and Luc Haine, \emph{Bispectral {D}arboux
  transformations: an extension of the {K}rall polynomials}, Int. Math. Res.
  Not. (1997), no.~8, 359--392.

\bibitem{GruenbaumYakimov02}
F.~Alberto Gr\"unbaum and Milen Yakimov, \emph{Discrete bispectral {D}arboux
  transformations from {J}acobi operators}, Pacific J. Math. \textbf{204}
  (2002), no.~2, 395--431.

\bibitem{Jha2026}
R.~K. Jha, \emph{A finite irrationality measure for {C}atalan's constant},
  Preprint v1, October 8, 2026,
  \href{https://doi.org/10.5281/zenodo.23237269}{doi:10.5281/zenodo.23237269};
  Lean sources at \url{https://github.com/rohitkrjha/catalan-measure}, 2026.

\bibitem{Arb}
Fredrik Johansson, \emph{Arb: efficient arbitrary-precision midpoint-radius
  interval arithmetic}, IEEE Trans. Comput. \textbf{66} (2017), no.~8,
  1281--1292.

\bibitem{KLS}
Roelof Koekoek, Peter~A. Lesky, and Ren\'e~F. Swarttouw, \emph{Hypergeometric
  orthogonal polynomials and their $q$-analogues}, Springer Monographs in
  Mathematics, Springer, Berlin, 2010.

\bibitem{Kratt99}
Christian Krattenthaler, \emph{Advanced determinant calculus}, S\'em. Lothar.
  Combin. \textbf{42} (1999), Art. B42q, 67 pp.

\bibitem{Kratt05}
\bysame, \emph{Advanced determinant calculus: a complement}, Linear Algebra
  Appl. \textbf{411} (2005), 68--166.

\bibitem{KrattRivoal08}
Christian Krattenthaler and Tanguy Rivoal, \emph{On a linear form for
  {C}atalan's constant}, South East Asian J. Math. Math. Sci. \textbf{6}
  (2008), no.~2, 3--15.

\bibitem{KrattZudilin19}
Christian Krattenthaler and Wadim Zudilin, \emph{Hypergeometry inspired by
  irrationality questions}, Kyushu J. Math. \textbf{73} (2019), no.~1,
  189--203.

\bibitem{Lai2026}
Li~Lai, \emph{The linear independence of $1$, $\zeta(3)$, and $\zeta(4)$},
  2026, Preprint, 5 October 2026 (v2: 15 October 2026).

\bibitem{Litsey2026}
Sean Litsey, \emph{An irrationality measure for {C}atalan's constant}, Hexagon
  preprint hexagon:2610.00187v4, \url{https://hexagonmath.org/2610.00187v4},
  2026.

\bibitem{LiuJiangZhang26}
Jingwen Liu, Kai Jiang, and Pingwen Zhang, \emph{Irrationality exponents of
  logarithms of positive rational numbers}, 2026, Preprint, 7 October 2026.

\bibitem{LZZ25}
Junqi Liu, Jujian Zhang, and Lihong Zhi, \emph{A formal proof of the
  irrationality of $\zeta(3)$ in {L}ean 4}, 2025.

\bibitem{Loewner34}
Karl L\"owner, \emph{\"{U}ber monotone {M}atrixfunktionen}, Math. Z.
  \textbf{38} (1934), 177--216.

\bibitem{Marcovecchio09}
Raffaele Marcovecchio, \emph{The {R}hin--{V}iola method for $\log 2$}, Acta
  Arith. \textbf{139} (2009), no.~2, 147--184.

\bibitem{Nesterenko85}
Yu.~V. Nesterenko, \emph{On the linear independence of numbers}, Vestnik
  Moskov. Univ. Ser. I Mat. Mekh. (1985), no.~1, 46--49, In Russian; English
  translation: Moscow Univ. Math. Bull. 40 (1985), no. 1, 69--74.

\bibitem{NikishinSorokin91}
E.~M. Nikishin and V.~N. Sorokin, \emph{Rational approximations and
  orthogonality}, Translations of Mathematical Monographs, vol.~92, American
  Mathematical Society, Providence, RI, 1991.

\bibitem{OpenAI2026}
{OpenAI}, \emph{Catalan's constant is irrational}, Preprint (September 24,
  2026), with a Lean formalization, available at
  \href{https://github.com/openai/math/blob/fd4aeeb2ee4fc729c18d98444fed42fd0529eeeb/preprints/Catalans-constant-is-irrational-September-24-2026/paper.pdf}{the
  original paper PDF} and
  \href{https://github.com/openai/math/tree/fd4aeeb2ee4fc729c18d98444fed42fd0529eeeb/lean/OAI/NumberTheory/Catalan}{the
  Lean formalization} (repository commit \texttt{fd4aeeb}, October 7, 2026),
  2026.

\bibitem{OpenAIPi2026}
\bysame, \emph{The irrationality exponent of $\pi$ is 2}, Preprint (September
  24, 2026),
  \url{https://github.com/openai/math/tree/main/preprints/The-irrationality-exponent-of-pi-is-2-September-24-2026},
  2026.

\bibitem{RhinViola96}
Georges Rhin and Carlo Viola, \emph{On a permutation group related to
  $\zeta(2)$}, Acta Arith. \textbf{77} (1996), no.~1, 23--56.

\bibitem{RhinViola01}
\bysame, \emph{The group structure for $\zeta(3)$}, Acta Arith. \textbf{97}
  (2001), no.~3, 269--293.

\bibitem{RivoalZudilin2003}
Tanguy Rivoal and Wadim Zudilin, \emph{Diophantine properties of numbers
  related to {C}atalan's constant}, Math. Ann. \textbf{326} (2003), no.~4,
  705--721.

\bibitem{SaffTotik}
Edward~B. Saff and Vilmos Totik, \emph{Logarithmic potentials with external
  fields}, Grundlehren der mathematischen Wissenschaften, vol. 316, Springer,
  Berlin, 1997.

\bibitem{SpiridonovZhedanov97}
Vyacheslav Spiridonov and Alexei Zhedanov, \emph{Spectral transformations,
  self-similar reductions and orthogonal polynomials}, J. Phys. A \textbf{30}
  (1997), no.~21, 7621--7637.

\bibitem{Sun2026}
Zhi-Wei Sun, \emph{Catalan's constant is irrational}, 2026.

\bibitem{Mathlib}
{The mathlib Community}, \emph{The {L}ean mathematical library}, Proceedings of
  the 9th ACM SIGPLAN International Conference on Certified Programs and Proofs
  (CPP 2020), 2020, pp.~367--381.

\bibitem{PNTPlus}
{The PrimeNumberTheoremAnd contributors}, \emph{The {PNT+} project: a {L}ean
  formalization of the prime number theorem and related results},
  \url{https://github.com/AlexKontorovich/PrimeNumberTheoremAnd}, The vendored
  compatibility additions used here are recorded in the repository's
  \texttt{NOTICE} and \texttt{PORTING.md}.

\bibitem{Wachs2026}
Dane Wachs, \emph{A note on a recent claimed proof of the irrationality of
  {C}atalan's constant}, 2026.

\bibitem{ZeilbergerZudilin20}
Doron Zeilberger and Wadim Zudilin, \emph{The irrationality measure of $\pi$ is
  at most $7.103205334137\ldots$}, Mosc. J. Comb. Number Theory \textbf{9}
  (2020), no.~4, 407--419.

\bibitem{Zhedanov97}
Alexei Zhedanov, \emph{Rational spectral transformations and orthogonal
  polynomials}, J. Comput. Appl. Math. \textbf{85} (1997), no.~1, 67--86.

\bibitem{Zudilin02Cheb}
Wadim Zudilin, \emph{A few remarks on linear forms involving {C}atalan's
  constant}, Chebyshevskii Sb. \textbf{3} (2002), no.~2(4), 60--70, In Russian.

\bibitem{Zudilin2003}
\bysame, \emph{An {A}p\'ery-like difference equation for {C}atalan's constant},
  Electron. J. Combin. \textbf{10} (2003), Research Paper 14.

\bibitem{Zudilin}
\bysame, \emph{A determinantal approach to irrationality}, Constr. Approx.
  \textbf{45} (2017), no.~2, 301--310.

\bibitem{Zudilin19Survey}
\bysame, \emph{Arithmetic of {C}atalan's constant and its relatives}, Abh.
  Math. Semin. Univ. Hambg. \textbf{89} (2019), no.~1, 45--53.

\end{thebibliography}


\end{document}
